% latexmk -pdf -shell-escape -interaction=nonstopmode masReferencia.tex
% latexmk -pdf -shell-escape -halt-on-error masReferencia.tex
\documentclass[8pt, spanish, mexico, letterpaper]{extarticle}
\usepackage[T1]{fontenc}
\usepackage[utf8]{inputenc}
\PassOptionsToPackage{hyphens}{url}\usepackage{hyperref}
\usepackage[landscape, left=0.45cm, right=0.45cm, top=0.5cm, bottom=0.45cm, includehead, headheight=14pt]{geometry}
\usepackage[spanish, es-tabla]{babel}
\usepackage[table,xcdraw]{xcolor}
\usepackage[italicdiff]{physics}
\usepackage{textcomp}
\usepackage{amsthm, amssymb, amsmath, amsfonts, bbm, mathtools}
\usepackage{minted}
\usepackage{enumerate}
\usepackage[shortlabels]{enumitem}
\usepackage{titling}
\usepackage{lmodern} %optimiza algunas fuentes
\usepackage{float}
\usepackage{mathrsfs}
\usepackage{wasysym}
\usepackage{csquotes}
\usepackage{dsfont}
\usepackage{caption}
\decimalpoint
\usepackage{graphicx}
\renewcommand{\baselinestretch}{1}
\usepackage{environ}

\usepackage{tocloft}
\usepackage{multicol}
\usepackage[skip=0pt]{parskip}
\usepackage{fancyhdr}
\usepackage{titlesec}

%%%%%%%%%%%%%%%%%% tabla de contenidos %%%%%%%%%%%%%%%%%%
\setlength{\cftbeforesecskip}{1pt} % Espacio entre secciones en el índice
\renewcommand{\cftsecdotsep}{\cftdotsep} % Poner puntitos a las secciones principales
\renewcommand{\cftsecfont}{\normalfont\footnotesize} % Letra más chica
\renewcommand{\cftsecpagefont}{\normalfont\footnotesize}
\renewcommand{\cftsubsecfont}{\normalfont\scriptsize}
\renewcommand{\cftsubsecpagefont}{\normalfont\scriptsize}

\newcommand{\Pa}[1]{\left( #1 \right)}
\newcommand{\PA}[1]{\big( #1 \big)}

%Images path
\graphicspath{{img/}}

%%%%%%%%%%%%%%%%%% d EXTENDED VERSION %%%%%%%%%%%%%%%%%%
\usepackage{etoolbox}
\newif\ifextendedversion
\extendedversiontrue   % Descomenta para la versión EXTENDIDA (Final Nacional)
% \extendedversionfalse  % Descomenta para la versión COMPACTA (Hiperregional)

\newcommand{\ext}[1]{\ifextendedversion#1\fi}

% 3. Bloque largo / Entorno (múltiples párrafos, listas, problemas tipo)
\NewEnviron{extended}{
    \ifextendedversion
        \vspace{0.3em}
        \begingroup
        \small\color{black!80}
        \vspace{0.2em}
        \noindent\textit{\footnotesize Versión Extendida:} \vspace{-1mm}\par
        \BODY
        \vspace{0.2em}
        \endgroup
        \vspace{0.3em}
    \fi
}


%%%%%%%%%%%%%%%%%% Minted %%%%%%%%%%%%%%%%%%
% \fontsize{tamaño de letra}{espacio entre líneas}\selectfont
\newcommand{\codesize}{\fontsize{7pt}{8pt}\selectfont} % tamfuente e interlineado para 9pt en el doc con extarticle
\setminted{
    fontsize=\codesize,      % Usa \scriptsize o \tiny si tienes muy buena vista
    baselinestretch=0.85,      % Comprime el espacio vertical entre líneas
    linenos=true,              % Números de línea (esenciales para debuggear en papel)
    numbersep=3pt,             % Pega los números de línea al código
    xleftmargin=12pt,          % Reduce el margen izquierdo que minted agrega por defecto
    breaklines=true,           % Ajuste de línea automático si te pasas del ancho
    breakanywhere=true,        % Rompe la línea incluso si no hay espacios
    tabsize=2,                 % Indentación corta (vital en 3 columnas)
    style=vs,                  % Tema visual claro y de alto contraste al imprimir
    stripall=true              % Elimina líneas en blanco al inicio y al final del archivo
}

% \inputcode{lenguaje}{ruta}{nombre_a_mostrar}
\newcommand{\inputcode}[3]{
    {\raggedleft \scriptsize \ttfamily \detokenize{#3} \par} 
    \vspace{-4pt}
    \inputminted[firstline=1, firstnumber=1]{#1}{#2}
}

%%%%%%%%%%%%%%%%%% Explicacion de codigo %%%%%%%%%%%%%%%%%%

% \algoinfo{Tiempo}{Memoria}{Explicación rápida}
\newcommand{\algoinfo}[3]{
    % \vspace{-2pt}
    \vspace{-1pt}
    {\fontsize{8pt}{9pt}\selectfont \textbf{T:} $O(#1)$ $\mid$ \textbf{M:} $O(#2)$ $\mid$ \textit{#3} \par}
    % \vspace{2pt}
}

%%%%%%%%%%%%%%%%%% Numeraciones %%%%%%%%%%%%%%%%%%
% Figuras numeradas
\usepackage{chngcntr}
\counterwithin{figure}{section}

% Ecuaciones numeradas
\numberwithin{equation}{section}

\allowdisplaybreaks

%%%%%%%%%%%%%%%%%% Definiciones de mate %%%%%%%%%%%%%%%%%%
\renewcommand{\theenumi}{\alph{enumi}}
\renewcommand{\labelenumi}{{\theenumi})}
\renewcommand\qedsymbol{$\blacksquare$}
\newcommand{\N}{\mathbb{N}}
\newcommand{\Z}{\mathbb{Z}}
\newcommand{\R}{\mathbb{R}}
\newcommand{\I}{\mathbb{I}}
\newcommand{\Q}{\mathbb{Q}}
\newcommand{\Zm}[1]{\mathbb{Z}/ #1 \mathbb{Z}}
\newcommand\scalemath[2]{\scalebox{#1}{\mbox{\ensuremath{\displaystyle #2}}}}

%%%%%%%%%%%%%%%%%% Comandos especiales de mate %%%%%%%%%%%%%%%%%%
\newcommand{\pa}[1]{\left( #1 \right)}
\newcommand{\br}[1]{\left[ #1 \right]}
\DeclarePairedDelimiter{\set}{\{}{\}}
\DeclarePairedDelimiter{\floor}\lfloor\rfloor
\DeclarePairedDelimiter{\ceil}\lceil\rceil
\DeclareMathOperator{\sgn}{sgn}
\DeclareMathOperator{\ord}{ord}
\newenvironment{solution}
{\renewcommand\xqedsymbol{$\square$}\begin{proof}[Solución]}
    {\end{proof}}

%%%%%%%%%%%%%%%%%% Espaciado %%%%%%%%%%%%%%%%%%
\setlength{\columnsep}{0.25cm} % Espacio minimo entre columnas
\setlength{\columnseprule}{0.5pt} % Linea divisoria opcional para guiar la vista
\setlength{\parindent}{0pt}
\setlist{topsep=0pt, partopsep=0pt, parsep=0pt, itemsep=0.5pt, leftmargin=3mm}
\setlist[enumerate]{label=\arabic*., leftmargin=3mm}
\setlist[itemize]{leftmargin=2.5mm}
\setlist[description]{leftmargin=0mm}


% \titlespacing*{comando}{izq}{arriba}{abajo}
% Espaciado hiper-compacto
\titlespacing*{\section}{0pt}{5pt}{3pt}
\titlespacing*{\subsection}{0pt}{3pt}{1pt}
\titlespacing*{\subsubsection}{0pt}{2pt}{1pt}

% Jerarquía visual (Pensada para impresión monocromática de ICPC)

% SECTION: MAYÚSCULAS y una línea negra sólida por debajo
\titleformat{\section}{\normalfont\normalsize\bfseries}{}{0pt}{\MakeUppercase}[{\vspace{1pt}\titlerule[0.8pt]}]

% SUBSECTION: Flecha indicadora, negrita (sin línea)
\titleformat{\subsection}{\normalfont\normalsize\bfseries}{}{0pt}{$\blacktriangleright$\hspace{3pt}}

% SUBSUBSECTION: Cursiva negrita y punto discreto (letra un poco más chica)
\titleformat{\subsubsection}{\normalfont\small\bfseries\itshape}{}{0pt}{$\bullet$\hspace{3pt}}

% Eliminar el salto de línea post-titulo (Run-in titles) si quieres compresión extrema:
% \titlespacing*{\subsection}{0pt}{1pt}{1em} % El tercer valor en 'em' hace que el texto siga en la misma línea

% Header
\pagestyle{fancy}
\headheight 14pt % Reducido de 25pt
\fancyhead{}
\fancyfoot{}
\lhead{\textbf{Mas Reference - ICPC} $\mid$ Mäthgic Crüe $\mid$ UG-CIMAT $\mid$ \today} % Todo en una línea
\rhead{ \textbf{Version Extra }- \textbf{Pág. \thepage}}
\renewcommand{\headrulewidth}{0.4pt}
\setlength\headsep{4pt} % Pegar el texto principal al header

% Comprimir espacio vertical alrededor de ecuaciones
\AtBeginDocument{
  \setlength{\abovedisplayskip}{2pt}
  \setlength{\belowdisplayskip}{2pt}
  \setlength{\abovedisplayshortskip}{0pt}
  \setlength{\belowdisplayshortskip}{0pt}
}

\usepackage[spanish, es-tabla]{babel}
\usepackage[table,xcdraw]{xcolor}
\newcommand{\E}{\mathbb{E}}
\newcommand{\Var}{\operatorname{Var}}
\newcommand{\Prob}{\mathbb{P}}

\newcommand{\rojo}[1]{\textcolor{red!80!black}{#1}}
\newcommand{\naranja}[1]{\textcolor{orange!50!red!80!black}{#1}}
\newcommand{\p}[2]{#1 \times 10^{#2}}

\newcommand{\<}{<}
% ---------- macros ----------
\newcommand{\sect}[1]{\medskip{\large\textbf{$\blacktriangleright$ #1}}\par\nopagebreak\smallskip}
\newcommand{\ov}[1]{\overline{#1}}
\newcommand{\sd}{\mathbin{\triangle}}
\newcommand{\Pw}{\mathcal{P}}
\newcommand{\ind}[1]{\mathbf{1}_{#1}}
\newcommand{\card}[1]{\left|#1\right|}
\newcommand{\bor}{\texttt{|}}
\newcommand{\band}{\texttt{\&}}
\DeclareMathOperator{\rad}{rad}
\DeclareMathOperator{\popcount}{popcount}

\begin{document}
\begin{multicols*}{3}
\raggedcolumns

\setlist[description]{
    % style=multiline,      % Mantiene un buen alineamiento si las etiquetas son largas
    leftmargin=02mm,     % Sangría uniforme para la descripción
    % labelwidth=0.0cm,     % Espacio reservado para el término (etiqueta)
    labelsep=01mm,       % Espacio entre el término y el texto
    itemsep=2pt,          % Espacio vertical entre ítems
    parsep=0pt,           % Espacio adicional entre párrafos dentro de un ítem
    topsep=0pt            % Espacio antes y después de toda la lista
}

\section{Plantilla y entorno}

\begin{minted}{cpp}
#include<bits/stdc++.h>
using namespace std;
// --- --- --- --- Type Aliases --- --- --- ---
typedef int64_t ll;
typedef pair<int, int> pii;
typedef pair<ll, ll> pll;
typedef vector<int> vi;
typedef vector<ll> vll;
typedef vector<vi> vvi;
typedef vector<vll> vvll;
typedef vector<pii> vpii;
typedef vector<pll> vpll;
// --- --- --- --- Short Macros --- --- --- ---
#define fi first
#define se second
#define pb push_back
#define all(x) (x).begin(), (x).end()
#define sz(x) (int)(x).size()
#define rep(i,a,b) for( int i=a; i< (b); ++i )
const int MOD = 1e9 + 7;
// --- --- --- --- LOGGER SETUP --- --- --- ---    
template<typename C, typename T = typename enable_if<!is_same<C, string>::value, typename C::value_type>::type>
ostream& operator<<(ostream &os, const C &v);
// Overload for std::pair
template<typename A, typename B>
ostream& operator<<(ostream &os, const pair<A, B> &p) {
    return os << '(' << p.fi << ", " << p.se << ')';
}
// Overload for Containers (excluding std::string) SFINAE
template<typename C, typename T>
ostream& operator<<(ostream &os, const C &v) {
    string sep;
    for(const T &x : v) os << sep << x, sep = " ";
    return os;
}
// colores
constexpr const char *RST="\033[0m", *GREEN="\033[1;32m", 
*YELLOW="\033[1;33m",*CYAN="\033[1;36m",*RED="\033[1;31m";
// Logger Function
template<typename ...Args>
void logger(const char* col, string vars, Args&&... values){
    cerr << col << '[' << vars << "] =" << RST;
    ((cerr << " [" << values << ']'), ...);
    cerr << '\n';
}
// --- --- --- --- DEBUGGER SETUP --- --- --- ---    
#ifdef LOCAL
  #define DEBUG   if (true)
  #define SALTO   (cerr << '\n')  
  #define LOG(...)   logger("", #__VA_ARGS__, __VA_ARGS__)
  #define LOGC(c, ...) logger(c, #__VA_ARGS__, __VA_ARGS__)    
  #define COL(c, ...) (cerr << c << __VA_ARGS__ << RST << '\n')
#else
    #define DEBUG   if (false)]
    #define SALTO
    #define LOG(...)
    #define LOGC(...)
    #define COL(...)
#endif
// // // // // // // // // // // // // // // // // // // // //
int main(){
    #ifdef LOCAL
        ifstream cin("in.in");
    #else
        ios_base::sync_with_stdio(0); cin.tie(0); cout.tie(0);
    #endif 
}
\end{minted}


En el archivo \texttt{\textasciitilde/.bashrc} (o \texttt{\textasciitilde/.zshrc} para macOS) escribir:
\begin{minted}{cpp}
alias gc='g++ -std=c++17 -O3 -Wall -Wextra -DLOCAL'
alias g='g++ -std=c++17 -O3 -Wall -Wextra'
source ~/.zshrc
\end{minted}

% \columnbreak
\subsection{Memoria disponible}

\begin{minted}{cpp} 
-Wl,-stack_size,0x8000000 // 128 MB  // 0x10000000 // 256 MB
-Wl,-stack_size,0x20000000 // 512 MB // 0x40000000 // 1 GB 
ulimit -s                 ulimit -s 524288 // 512MB 
ulimit -s unlimited       ulimit -s 1048576 // 1gb
\end{minted}



\section{Informacion adicional}

\subsection*{Constantes y más}

\subsubsection*{Límites de tipos}
\begin{tabular}{ll}
\texttt{INT\_MAX} & \(2{,}147{,}483{,}647 = 2^{31}-1 \approx 2.1\times10^{9}\)\\
\texttt{INT\_MIN} & \(-2^{31} = -2{,}147{,}483{,}648\)\\
\texttt{UINT\_MAX} & \(4{,}294{,}967{,}295 = 2^{32}-1 \approx 4.29\times10^{9}\)\\
\texttt{LLONG\_MAX} & \(9{,}223{,}372{,}036{,}854{,}775{,}807 = 2^{63}-1 \approx 9.22\times10^{18}\)\\
\texttt{LLONG\_MIN} & \(-2^{63}\)\\
\texttt{ULLONG\_MAX} & \(2^{64}-1 = 18{,}446{,}744{,}073{,}709{,}551{,}615 \approx 1.84\times10^{19}\)\\
\texttt{double} & enteros exactos hasta \(2^{53} \approx 9.007\times10^{15}\)\\
\end{tabular}

\subsubsection*{Valores de INF recomendados}
\texttt{int}: \(\texttt{INF}=10^9\) (o \texttt{0x3f3f3f3f} $\approx 1.06\times10^9$, permite $\texttt{INF}+\texttt{INF}$ sin overflow). \\
\texttt{long long}: \(\texttt{INF}=10^{18}\) (o \texttt{4e18}; ojo con sumar dos INF).

\subsubsection*{Potencias de 2}
\begin{tabular}{rl|rl|rl}
\(2^{10}\) & 1,024 & \(2^{19}\) & 524,288 & \(2^{37}\) & \(\approx 1.37\times10^{11}\) \\
\(2^{11}\) & 2,048 & \(2^{20}\) & 1,048,576 & \(2^{40}\) & \(\approx 1.10\times10^{12}\) \\
\(2^{12}\) & 4,096 & \(2^{22}\) & 4,194,304 & \(2^{45}\) & \(\approx 3.52\times10^{13}\) \\
\(2^{13}\) & 8,192 & \(2^{25}\) & 33,554,432 & \(2^{50}\) & \(\approx 1.13\times10^{15}\) \\
\(2^{14}\) & 16,384 & \(2^{27}\) & 134,217,728 & \(2^{60}\) & \(\approx 1.15\times10^{18}\) \\
\(2^{15}\) & 32,768 & \(2^{30}\) & \(\approx 1.07\times10^9\) & \(2^{63}\) & \(\approx 9.22\times10^{18}\) \\
\(2^{16}\) & 65,536 & \(2^{31}\) & \(\approx 2.15\times10^9\) & \(2^{64}\) & \(\approx 1.84\times10^{19}\) \\
\(2^{17}\) & 131,072 & \(2^{32}\) & \(\approx 4.29\times10^9\) & \(2^{127}\) & \(\approx 1.70\times10^{38}\) \\
\(2^{18}\) & 262,144 & \(2^{35}\) & \(\approx 3.44\times10^{10}\) & \(2^{128}\) & \(\approx 3.40\times10^{38}\) \\
\end{tabular}

\subsubsection*{ASCII}
\texttt{'0'}=48, \texttt{'A'}=65, \texttt{'a'}=97 \quad (\texttt{'a'-'A'}=32, es decir, el bit 5).
\subsubsection*{Constantes reales}
\(\pi \approx 3.14159265358979\),  \(e \approx 2.71828182845905\), \(\sqrt2 \approx 1.41421356237\) \\
\(\varphi=\frac{1+\sqrt5}{2}\approx 1.6180339887\), \  \(\ln2\approx0.693147\), \quad \(\gamma\approx0.5772156649\). \\
\texttt{acos(-1)} da $\pi$ con precisión de \texttt{double}.

\subsubsection*{Complejidad según $n$ (límite $\approx 10^8$ ops simples/s)}
\begin{tabular}{ll|ll}
    $n \le 10$ & \(O(n!)\) & $n \le 5{,}000$ & \(O(n^2)\) \\
    $n \le 20$--$25$ & \(O(2^n\cdot n)\) & $n \le 10^5$--$2\times10^5$ & \(O(n\log n)\) \\
    $n \le 100$ & \(O(n^4)\) & $n \le 10^6$--$10^7$ & \(O(n)\) \\
    $n \le 500$ & \(O(n^3)\) & $n \le 10^{18}$ & \(O(\log n)\) o \(O(1)\) \\
\end{tabular}

\subsubsection*{Factoriales y Fibonacci (máximo que cabe)}
\(12! = 479{,}001{,}600\) es el mayor en \texttt{int}; \(13!>\texttt{INT\_MAX}\).\\
\(20! = 2{,}432{,}902{,}008{,}176{,}640{,}000\) es el mayor en \texttt{long long}.\\
\(F_{46} = 1{,}836{,}311{,}903\) mayor en \texttt{int};\\
\(F_{92}=7{,}540{,}113{,}804{,}746{,}346{,}429\) mayor en \texttt{long long}.\\
\(10! = 3{,}628{,}800\), \\
Catalan:  \(C_{15}=9{,}694{,}845\), \(C_{19}=1{,}767{,}263{,}190\) (mayor en \texttt{int}).

\begin{description}
\item[Máximo número de divisores]
$n\le10^6$: 240 ($720720$) \quad $n\le 10^9$: 1344 ($735134400$) \quad $n\le10^{18}$: 103,680 ($897612484786617600$).

\item[Primos y módulos comunes]
\(10^9+7\), \quad \(998{,}244{,}353 = 119\cdot2^{23}+1\) (NTT), \quad \(10^9+9\).
\end{description}

\subsection{Primos (1--1000)(168)}
\tiny
2, 3, 5, 7, 11, 13, 17, 19, 23, 29, 31, 37, 41, 43, 47, 53, 59, 61, 67, 71, 73, 79, 83, 89, 97, 101, 103, 107, 109, 113, 127, 131, 137, 139, 149, 151, 157, 163, 167, 173, 179, 181, 191, 193, 197, 199, 211, 223, 227, 229, 233, 239, 241, 251, 257, 263, 269, 271, 277, 281, 283, 293, 307, 311, 313, 317, 331, 337, 347, 349, 353, 359, 367, 373, 379, 383, 389, 397, 401, 409, 419, 421, 431, 433, 439, 443, 449, 457, 461, 463, 467, 479, 487, 491, 499, 503, 509, 521, 523, 541, 547, 557, 563, 569, 571, 577, 587, 593, 599, 601, 607, 613, 617, 619, 631, 641, 643, 647, 653, 659, 661, 673, 677, 683, 691, 701, 709, 719, 727, 733, 739, 743, 751, 757, 761, 769, 773, 787, 797, 809, 811, 821, 823, 827, 829, 839, 853, 857, 859, 863, 877, 881, 883, 887, 907, 911, 919, 929, 937, 941, 947, 953, 967, 971, 977, 983, 991, 997, 

\normalsize

\columnbreak
\subsubsection*{Cantidades de primos $\le n$:}
\begin{center}

\begin{tabular}{cc | cc }
$n$ & $\pi(n)$ & $n$ & $\pi(n)$ \\ \hline
$10^2$ & 25 & $10^9$ & 50,847,534  \\
$10^3$ & 168 & $10^{10}$ & 455,052,511  \\
$10^4$ & 1,229 & $10^{11}$ & 4,118,054,813 \\
$10^5$ & 9,592 & $10^{12}$ & 37,607,912,018  \\
$10^6$ & 78,498 & $10^{13}$ & 346,065,536,839  \\
$10^7$ & 664,579 & $10^{14}$ & 3,204,941,750,802  \\
$10^8$ & 5,761,455 & $10^{15}$ & 29,844,570,422,669  
\end{tabular}
\end{center}








\subsection*{Aritmética entera}





\begin{itemize}
\item Para enteros \textbf{positivos} \texttt{int a,b,x} (o \texttt{long long}):
$$ \texttt{if( a*b > x )} \equiv \texttt{if( a > x/b )}$$
$$\texttt{if( a*b >= x )}\equiv
\texttt{if(a >= (x+b-1) / b)}
    \equiv \texttt{if(a > (x-1) / b)} $$
  \item Techo (positivos): con $a\ge1$
  $$ \lceil a/b\rceil = \texttt{(a + b - 1) / b} = \texttt{(a - 1) / b + 1}$$ 
  \item Módulo no negativo: \texttt{((a \% m) + m) \% m}. En C++ \texttt{-7 / 2 = -3} (trunca hacia 0) y \texttt{-7 \% 2 = -1}.
  \item Resta modular: \texttt{(a - b + m) \% m}. 
  \item Inverso modular ($m$ primo): \(a^{-1}\equiv a^{m-2}\pmod m\).
  \item \texttt{lcm(a,b) = a / gcd(a,b) * b} (div evitar overflow).
  \item Overflow de suma: \texttt{a + b > MAX} $\equiv$ \texttt{a > MAX - b}.
  \item Overflow de producto: equiv previa o \texttt{\_\_int128}.
  \item \texttt{(a + b) / 2} puede desbordar; usar \texttt{a + (b - a) / 2}.
\end{itemize}


\subsection*{Memoria}

\paragraph{Unidades.}
\(1\,\text{KB} = 2^{10} \) B, \quad \(1\,\text{MB} = 2^{20} =\) B, \quad \(1\,\text{GB} = 2^{30}\) B.\\
\textbf{Aprox:} \(1\,\text{MB} \approx 10^{6}\) B 
$2^{18}$ ints $=1$ MB, \  $2^{20}$ ints $=4$ MB.


\textbf{Tamaño de tipos (x86-64, GCC).} \\
\begin{tabular}{ll|ll}
\texttt{char}, \texttt{bool} & 1 B & \texttt{long long}, \texttt{double} & 8 B \\
\texttt{short} & 2 B & \texttt{long double}, \texttt{\_\_int128} & 16 B \\
\texttt{int}, \texttt{float} & 4 B & puntero (\texttt{int*}) & 8 B \\
\texttt{pair<int,int>} & 8 B & \texttt{pair<ll,ll>} & 16 B \\
\texttt{array<int,3>} & 12 B & \texttt{array<ll,3>} & 24 B 
\end{tabular}
\textbf{Aprox:} En la práctica con 256 MB: $\approx 5\times10^{7}$ ints o $\approx 2.5\times10^{7}$ long long.\\
\textbf{De $n$ elementos a MB.}\\
\begin{tabular}{r|rrrr}
$n$ & 1 B (\texttt{char}) & 4 B (\texttt{int}) & 8 B (\texttt{ll}) & 16 B \\ \hline
$10^5$ & 0.1 MB & 0.4 MB & 0.8 MB & 1.5 MB \\
$10^6$ & 0.95 MB & 3.8 MB & 7.6 MB & 15.3 MB \\
$10^7$ & 9.5 MB & 38.1 MB & 76.3 MB & 152.6 MB \\
$10^8$ & 95.4 MB & 381.5 MB & 762.9 MB & 1,526 MB 
\end{tabular}


\textbf{De MB a $n$ elementos.}\\
\begin{tabular}{r|rrrr}
Límite & \texttt{char} (1 B) & \texttt{int} (4 B) & \texttt{ll} (8 B) & 16 B \\ \hline
64 MB  & $6.7\times10^{7}$ & $1.68\times10^{7}$ & $8.4\times10^{6}$ & $4.2\times10^{6}$ \\
128 MB & $1.34\times10^{8}$ & $3.36\times10^{7}$ & $1.68\times10^{7}$ & $8.4\times10^{6}$ \\
256 MB & $2.68\times10^{8}$ & $6.71\times10^{7}$ & $3.36\times10^{7}$ & $1.68\times10^{7}$ \\
512 MB & $5.37\times10^{8}$ & $1.34\times10^{8}$ & $6.71\times10^{7}$ & $3.36\times10^{7}$ \\
1024 MB & $1.07\times10^{9}$ & $2.68\times10^{8}$ & $1.34\times10^{8}$ & $6.71\times10^{7}$ \\
\end{tabular}\\



% \paragraph{Bitset y criba.}
% \texttt{bitset<N>} usa $N/8$ bytes: \(10^8\) bits $= 12.5$ MB, \(10^9\) bits $\approx 119$ MB.\\
% Criba de Eratóstenes: \texttt{bool} hasta $10^7$ $\approx 9.5$ MB, hasta $10^8$ $\approx 95$ MB; con \texttt{bitset} hasta $10^9$ $\approx 119$ MB.

% \paragraph{Estructuras con tamaño típico.}
% \begin{itemize}
%   \item \texttt{vector<int>} de $n$ elementos: $4n$ B + 24 B de cabecera. Al crecer con \texttt{push\_back} la capacidad se duplica, así que el pico puede ser $\approx 2\times$ (usar \texttt{reserve} si se conoce $n$).
%   \item \texttt{vector<vector<int>>} (grafo): 24 B por vértice + 4 B por cada extremo de arista. Para $m$ aristas no dirigidas: $\approx 8m$ B de datos.
%   \item Lista de aristas \texttt{(u,v,w)} de \texttt{int}: 12 B por arista. Con $m=2\times10^5$ son $\approx 2.3$ MB.
%   \item Árbol de segmentos de tamaño $4n$: \(n=10^5\): 1.5 MB (\texttt{int}) / 3.1 MB (\texttt{ll}); \quad \(n=10^6\): 15.3 MB / 30.5 MB.
%   \item Fenwick / prefijos: $n$ elementos, igual que un arreglo normal.
%   \item Tabla sparse: $n\log_2 n$ elementos. \(n=10^5\): $\approx 17\times10^5$ ints $=6.5$ MB; \quad \(n=10^6\): $\approx 20\times10^6$ ints $=76$ MB.
%   \item \texttt{set<int>} / \texttt{map<int,int>}: $\approx 40$--$48$ B por elemento (nodo + punteros). $10^6$ elementos $\approx 45$ MB. Mucho más que un arreglo.
%   \item \texttt{unordered\_map}: $\approx 32$--$40$ B por elemento más la tabla de buckets.
%   \item \texttt{string}: 32 B de objeto (los strings cortos caben dentro); un string largo ocupa su longitud + 32 B.
% \end{itemize}

% \paragraph{Stack y recursión.}
% Cada llamada recursiva gasta $\approx 32$--$100$ B. Una DFS de profundidad $10^6$ puede usar $\approx 50$--$100$ MB de stack. En Codeforces el stack es grande, pero en otros jueces (o con el valor por defecto de Linux, 8 MB) puede dar \emph{stack overflow}. Para profundidades de $10^5$--$10^6$ conviene una versión iterativa.

% \paragraph{Reglas rápidas.}
% \begin{itemize}
%   \item Arreglo global de $10^6$ \texttt{int} $\approx 4$ MB: no hay problema.
%   \item Contar siempre \emph{todos} los arreglos: 20 arreglos de $10^6$ \texttt{ll} son 160 MB.
%   \item Si $n\cdot s$ pasa el límite, pensar en \texttt{bitset}, \texttt{short}/\texttt{char}, comprimir coordenadas, o DP con filas rodantes.
%   \item Los arreglos globales no usan stack; los locales grandes (\texttt{int a[10000000]} dentro de \texttt{main}) sí pueden reventarlo.
%   \item Un \texttt{Memory Limit Exceeded} con pocos datos suele indicar un \texttt{vector} que crece sin control o un \texttt{map} muy grande.
% \end{itemize}



%%%%%%
\columnbreak

\section{Ideas  < - < - < - }

\subsection{DP}
\begin{description}
\item[Elegir el estado]
   No siempre pienses en prefijos $[0,i]$; considera también sufijos.
    \begin{itemize}
      \item $dp[i][\dots]$ = es posible sobrevivir las primeras $i$ horas, $[1,i]$.
      \item $dp[i][\dots]$ = es posible sobrevivir desde el fin de la hora $i$ hasta el final, $(i,n+1]$.
    \end{itemize}


\item[Transiciones sobre el último elemento]
  Revisar si ...
  \begin{itemize}
    \item se usa el último elemento o no.
    \item se usa el último elemento un número par o impar de veces.
  \end{itemize}

\item[Distinguir entre dos formulaciones]
  $dp[\dots][i]$ :
  \begin{itemize}
    \item $dp[\dots][i]$ = ya completé los primeros $i$ elementos.
    \item $dp[\dots][i]$ = el $i$-ésimo elemento queda completado/terminado.
  \end{itemize}

\item[Probar quien es el ultimo elemento]
  Si $n \le 20$ y el orden importa, la transición desde $dp[mask]$ decide qué elemento se agrega \emph{al final}.
\end{description}


\subsection{Búsqueda binaria}
\begin{description}
\item[Búsqueda binaria sobre la respuesta]
  Lograr $S$ en a lo sumo $K$ pasos: binaria sobre $S$. Para $S$ fijo, calcular $f(S)$, la cantidad mínima de pasos para obtener $S$, y buscar $\max\{S : f(S) \le K\}$.
  \textbf{Requisito:} $f$ debe ser monótona (si se puede lograr $S$, se puede lograr cualquier $S' < S$).
\end{description}

\subsection{Ideas más generales}
\begin{description}
    \item[Ciclos] Avanzar un paso a la vez, no varios.
    \item[Local] Pasar de condiciones locales a condiciones globales.
    \item[Independencia]  (no probabilística): decisiones o eventos que no se afectan entre sí, por ejemplo poner G y T.
    \item[Bases o generadores:]  \textcolor{white}{.}
      \begin{itemize}
        \item una base $B$ de un espacio vectorial;
        \item los generadores de una $\sigma$-álgebra, $\sigma(\mathcal{C})$.
      \end{itemize}
    \item[Informacion] En el estado actual, ¿qué información se puede olvidar y cuál es necesario conservar para resolver el problema?
\end{description}

\subsection{Ideas especificas}
\begin{description}
    \item[No buscar todo] Los problemas ``\textbf{South America} 2025 - Brazil First Phase \textbf{K} - Knockout'' y ``(\textbf{NWERC} 2025) \textbf{D} Dreamcatcher'' mostraron que no hace falta probar toda la fuerza bruta si no algunos ``lugares cercanos a donde parece que debería estar el maximo''. En lugar de probar $O(N)$ probar una \textbf{cierta cantidad}, o empezar a probar \textbf{desde ``el lugar donde este el maximo'' }
\end{description}

\section{Observaciones al implementar}
\begin{description}
  \item[Productos en \texttt{int}] Castea \emph{antes} de multiplicar: \texttt{(ll)(i-1)*i/2}. Castear el resultado ya no sirve.
  \item[Saturación] Si un valor se acota a \texttt{INF}, trátalo como ``$\geq$ INF'', nunca como 0. Guarda el valor saturado \emph{antes} de cualquier \texttt{break}.
  \item[Cortes en DP] Cortar el cálculo solo es válido si lo que no se calcula nunca se lee después. Verifica que los índices leídos siempre queden dentro del rango calculado.
  \item[Restas con saturación] Una resta sobre valores saturados es inexacta. Asegura que el sustraendo siempre sea exacto, o usa \texttt{\_\_int128} para los intermedios.
  \item[Intermedios anchos] Calcula en \texttt{\_\_int128} y satura al guardar. La suma de dos valores $\leq 4\cdot10^{18}$ cabe en \texttt{ll}, pero una suma con resta intermedia puede no caber.
  \item[Monotonía como test] Si sabes que una cantidad es monótona, ponle un \texttt{assert}. Un salto raro suele indicar un bug de truncamiento y no de overflow.
  \item[Depuración] Imprime en puntos de transición (donde cambia el comportamiento) y no en intervalos regulares.
  \item[Memoria] Si la recurrencia solo usa la fila anterior, usa dos filas. Una tabla de \texttt{ll} de $2\cdot10^5\times196$ pesa unos 313 MB.
\end{description}


\section{Cosas triviales a tomar en cuenta}
\subsection*{  ¿Qué hacer si da \textbf{\rojo{WA} (\rojo Wrong \rojo Answer)}? }
\begin{itemize}
    \item Revisar que se esta impriminedo. Otros couts, formato.
    \item Quitar \texttt{ifstream cin(``in.txt'');} y revisar \texttt{DEBUG}
    \item \textbf{Revisar Overflow:} revisar los tipos de datos (\texttt{int}, \texttt{int64\_t}).
    \item \textbf{Revisar línea a línea:} asegúrate de que cada linea si este correcta, por ejemplo \texttt{fact[n]} y no solo \texttt{n}.
    \item \textbf{Casos extremos:} prueba con casos límites, como $N = 1$, o valores muy grandes.
    \item  Los \text{for}s estan bien: no se sale o empieza antes. Por ejemplo en divisibildiad y primos tener cuidado si se inicia en \texttt{i=1} en lugar de \texttt{i=2}
    \item  Revisar los limites del problema y los limites que se pusieron en el codigo tipo \texttt{MAXN= }, \texttt{r= }, etc.
    \item  Revisar como llenas la dp, o como llenas arreglos, sobretodo si es inplace, no es lo mismo $0\to n$ que $n\to 0$.
    \item  ... 
\end{itemize}

\subsection*{¿Qué hacer si da \textbf{\naranja{RT} (\naranja Runtime  \naranja error)}?}
\begin{itemize}
    \item  Declarar arreglos fuera del main.
    \item  Si usas modulos, ojo con las restas aunque creas que da postivio. $A>B$ \textbf{NO} implica $(A \mod m )> (B \mod m)$.
    \item  ... 
\end{itemize}


\subsection*{¿Qué hacer si da \textbf{\naranja{TLE} (\naranja Time \naranja Limit \naranja Excede)}? }
\begin{itemize}
    \item Quitar el Print Debugger 
    \item   Revisar que si este usando
    \begin{minted}{cpp}
    ios_base::sync_with_stdio(0);
    cin.tie(0);
    cout.tie(0);
    \end{minted}
    \item Ver que nada se esta ciclando
    \item 
\end{itemize}


\subsection*{¿Qué hacer si da \textbf{\naranja{MLE} (\naranja Memory   \naranja Limit \naranja  Exceeded) }? }
Para medir el consumo pico de memoria RAM
\begin{minted}{bash}
g++ -O3 j.cpp -o j
/usr/bin/time -v ./j < jenga.in
\end{minted}

En la salida generada, busca la línea:
\begin{center}
\texttt{Maximum resident set size (kbytes): [valor]}
\end{center}
Divide el valor obtenido entre $1024$ para convertirlo a megabytes (MB) y compararlo con el límite del problema.

\subsection*{Memoria}

\begin{description}
    \item[\texttt{T arr[N]}] $N \times \text{sz(T)}$ -- La más eficiente (contigua).
    \item[\texttt{vector<T>}] $N \times \text{sz(T)}$ -- +32 bytes de overhead del objeto.
    \item[\texttt{vector<bool>}] $\approx N / 8$ bytes -- Optimizado (1 bit por elemento).
    \item[\texttt{bitset<N>}] $\approx N / 8$ bytes -- Estático, el más ligero de todos.
    \item[\texttt{set<T>}] $\approx 48N$ bytes -- Árbol: 3 punteros + color + dato + padding.
    \item[\texttt{map<K, V>}] $\approx 48N$ bytes -- Similar al set, guarda pares \texttt{std::pair}.
    \item[\texttt{unordered\_set<T>}] $\approx 32N$ bytes (Tabla Hash: ptr por nodo + overhead de buckets)
    \item[\texttt{list<T>}] $\approx 24N + \text{sz(T)} \times N$ -- List doblement enlazada (2 punteros).
    \item[\texttt{deque<T>}] $\approx N \times \text{sz(T)}$(Fragment en bloques,más pesada que vector)
    \item[\texttt{priority\_queue<T>}] $N \times \text{sz(T)}$ -- Usa un \texttt{vector} por debajo (Heap).
    \item[\texttt{set<T>}] $32 \text{ a } 48 \text{ bytes} + \text{sz(T)}$ (Memoria por nodo, aprox. 64-bit).
    \item[\texttt{map<K, V>}] $32 \text{ a } 48 \text{ bytes} + \text{sz(K)} + \text{sz(V)}$ 
    \item[\texttt{unordered\_set<T>}] $24 \text{ a } 32 \text{ bytes} + \text{sz(T)}$ 
\end{description}


\section{Sobrecargar}
\begin{minted}{cpp}
struct T {
    ll a, b;
    // ctor por defecto + parametros
    T(ll a_ = 0, ll b_ = 0) : a(a_), b(b_) {}          
    //  compuestos: miembro, modifican *this y devuelven T& 
    T& operator+=(const T& o) 
        { a += o.a; b += o.b; return *this; }
    T& operator-=(const T& o) 
        { a -= o.a; b -= o.b; return *this; }
    T& operator*=(ll k)  { a *= k;   b *= k;   return *this; }
    // --- binarios: friend, copia por valor + compuesto ---
    friend T operator+(T x, const T& y) { return x += y; }
    friend T operator-(T x, const T& y) { return x -= y; }
    friend T operator*(T x, ll k)       { return x *= k; }
    // k*x conmutado
    friend T operator*(ll k, T x)       { return x *= k; }  
    // unario
    T operator-() const { return T(-a, -b); }   
    // ++t  (prefijo, devuelve referencia)
    T& operator++()    { ++a; ++b; return *this; }   
    // t++ (postfijo, devuelve la copia vieja)
    T  operator++(int) { T r = *this; ++*this; return r; } 
    //  comparacion (set, map, sort, priority_queue usan < ) 
    bool operator<(const T& o) const 
        { return tie(a, b) < tie(o.a, o.b); }
    bool operator==(const T& o) const 
        { return a == o.a && b == o.b; }
    bool operator!=(const T& o) const 
        { return !(*this == o); }
    // C++20: reemplaza los 3 anteriores (y >, <=, >=) con:
    auto operator<=>(const T& o) const 
        { return tie(b, a) <=> tie(o.b, o.a); }
    bool operator==(const T& o) const = default;   
    // --- otros utiles ---
    // t(x) como funcion
    ll operator()(ll x) const { return a * x + b; }         
    // t[0], t[1]
    ll& operator[](int i) { return i ? b : a; }             
    // if(t)
    explicit operator bool() const { return a || b; }       
    // --- I/O: free functions, NO miembros ---
    friend ostream& operator<<(ostream& os, const T& t) 
        { return os << t.a << ' ' << t.b; }
    friend istream& operator>>(istream& is, T& t)       
        { return is >> t.a >> t.b; }
};
// hash para unordered_map/set 
// (fuera del struct, requiere operator==)
template<> struct std::hash<T> {
    size_t operator()(const T& t) const { 
        return hash<ll>()(t.a * 1000003 ^ t.b); 
    }
};
\end{minted}



%=%=%=%=%=%=%=%=%=%=%=%=%=%=%=%=%=%=%=%=%=%=%=%=%=%=
%=%=%=%=%=%=%=%=%=%=%=%=%=%=%=%=%=%=%=%=%=%=%=%=%=%=
%=%=%=%=%=%=%=%=%=%=%=%=%=%=%=%=%=%=%=%=%=%=%=%=%=%=
%=%=%=%=%=%=%=%=%=%=%=%=%=%=%=%=%=%=%=%=%=%=%=%=%=%=
\clearpage
\section{Teoría de Conjuntos}

Universo $U$, complemento $\ov{A}=U\setminus A$, diferencia sim\'etrica
$A\sd B=(A\setminus B)\cup(B\setminus A)$ (es el $\oplus$ de la hoja de l\'ogica).
Toda identidad se prueba \emph{elemento por elemento} o con indicadores (\S2).
 
% =====================================================
\subsection{Diccionario conjuntos -- l\'ogica -- bitwise}
\begin{center}
\begin{tabular}{@{}lll@{}}
\hline
Conjuntos & L\'ogica ($x\in\cdot$) & M\'ascara \\
\hline
$A\cup B$ & $p\lor q$ & \texttt{a\bor b} \\
$A\cap B$ & $p\land q$ & \texttt{a\band b} \\
$\ov A$ & $\lnot p$ & \texttt{\textasciitilde a \band{} full} \\
$A\setminus B$ & $p\land\lnot q$ & \texttt{a \band{} \textasciitilde b} \\
$A\sd B$ & $p\oplus q$ & \texttt{a\textasciicircum b} \\
$A\subseteq B$ & $p\Rightarrow q$ & \texttt{(a\band b)==a} \\
$\card{A}$ & --- & \texttt{popcount(a)} \\
\hline
\end{tabular}
\end{center}
 
% =====================================================
\subsection{\'Algebra de Boole de $\Pw(U)$}
\begin{itemize}
  \item $A\cup A=A\cap A=A$ \quad (idempotencia); conmutatividad y asociatividad de $\cup,\cap,\sd$.
  \item $A\cup\varnothing=A$, \; $A\cap U=A$, \; $A\cup U=U$, \; $A\cap\varnothing=\varnothing$.
  \item Absorci\'on: $A\cup(A\cap B)=A$, \quad $A\cap(A\cup B)=A$.
  \item Complemento: $A\cup\ov A=U$, \; $A\cap\ov A=\varnothing$, \; $\ov{\ov A}=A$.
  \item $A\cap(B\cup C)=(A\cap B)\cup(A\cap C)$.
  \item $A\cup(B\cap C)=(A\cup B)\cap(A\cup C)$.
  \item De Morgan: $\ov{A\cup B}=\ov A\cap\ov B$, \quad $\ov{A\cap B}=\ov A\cup\ov B$.
  \item $A\subseteq B \iff A\cup B=B \iff A\cap B=A \iff A\setminus B=\varnothing \iff \ov B\subseteq\ov A$.
\end{itemize}
 
Diferencia:
\begin{itemize}
  \item $A\setminus B=A\cap\ov B=A\setminus(A\cap B)$.
  \item $(A\setminus B)\setminus C=A\setminus(B\cup C)$.
  \item $A\setminus(B\setminus C)=(A\setminus B)\cup(A\cap C)$.
  \item $A\setminus(B\cup C)=(A\setminus B)\cap(A\setminus C)$.
  \item $A\setminus(B\cap C)=(A\setminus B)\cup(A\setminus C)$.
  \item $(A\cup B)\setminus C=(A\setminus C)\cup(B\setminus C)$.
  \item $A\cap(B\setminus C)=(A\cap B)\setminus C$.
\end{itemize}

 
% =====================================================
\subsection{Diferencia sim\'etrica}
\begin{itemize}
  \item $A\sd B=(A\cup B)\setminus(A\cap B)$.
  \item $A\sd\varnothing=A$, \; $A\sd A=\varnothing$; por tanto $A\sd B=C\iff A=B\sd C$.
  \item $A\cup B=(A\sd B)\sqcup(A\cap B)$ \quad 
  \item $A\sd(A\cap B)=A\setminus B=(A\cup B)\sd B$  
  \item $(A\cap B)\sd(A\cup B)=A\sd B$ 
  \item $\ov A\sd\ov B=A\sd B$, \quad $\ov{A\sd B}=\ov A\sd B$.
  \item $\cap$ distribuye sobre $\sd$: $A\cap(B\sd C)=(A\cap B)\sd(A\cap C)$
  \item $(\Pw(U),\sd,\cap)$ es un anillo booleano.
  \item $\cup$ \textbf{no} distribuye sobre $\sd$. Contraejemplo: $A=\{1\}$, $B=C=\{2\}$:
        $A\cup(B\sd C)=\{1\}$ pero $(A\cup B)\sd(A\cup C)=\varnothing$.
  \item $\card{A\sd C}\le\card{A\sd B}+\card{B\sd C}$ 
  \item distancia de Hamming: $\card{A\sd B}=\popcount(a\oplus b)$)
  \item $\bigtriangleup_{i=1}^n A_i=\{x:\ x \text{ est\'a en un n\'umero impar de los } A_i\}$.
\end{itemize}
 
% =====================================================
\subsection{Indicadores y cardinalidades}
Con $\ind{A}(x)\in\{0,1\}$ y $\card A=\sum_x\ind A(x)$:
\begin{itemize}
  \item $\ind{A\cap B}=\ind A\,\ind B$, \quad $\ind{\ov A}=1-\ind A$.
  \item $\ind{A\cup B}=\ind A+\ind B-\ind A\,\ind B$.
  \item $\ind{A\sd B}=\ind A+\ind B-2\,\ind A\,\ind B\equiv\ind A+\ind B\pmod 2$.
  \item $\card A+\card B=\card{A\cup B}+\card{A\cap B}$ \ (an\'alogo $a+b=a\bor b+a\band b$).
  \item $\card{A\sd B}=\card A+\card B-2\card{A\cap B}$ \  (an\'alogo  $a+b=a\oplus b+2(a\band b)$).
  \item $\card{A\setminus B}=\card A-\card{A\cap B}$.
  \item $\max(\card A,\card B)\le\card{A\cup B}\le\card A+\card B$.
  \item $\card{A\cap B}\ge\card A+\card B-\card U$.
\end{itemize}
 
% =====================================================
\subsection{Uniones e intersecciones generalizadas}
\begin{itemize}
  \item $\bigcup_{i\in I}A_i=\{x:\exists i\in I,\ x\in A_i\}$, 
  \item $\bigcap_{i\in I}A_i=\{x\in U:\forall i\in I,\ x\in A_i\}$.
  \item Convenci\'on: $\bigcup_{\varnothing}=\varnothing$, \; $\bigcap_{\varnothing}=U$ 
  \item el $\exists$ sobre vac\'io es falso, el $\forall$ es verdadero.
  \item De Morgan: $\ov{\bigcup_iA_i}=\bigcap_i\ov{A_i}$, \quad $\ov{\bigcap_iA_i}=\bigcup_i\ov{A_i}$ \quad ($\lnot\exists\equiv\forall\lnot$).
  \item $B\cap\bigcup_iA_i=\bigcup_i(B\cap A_i)$, \quad $B\cup\bigcap_iA_i=\bigcap_i(B\cup A_i)$.
  \item $B\setminus\bigcup_iA_i=\bigcap_i(B\setminus A_i)$, \quad $B\setminus\bigcap_iA_i=\bigcup_i(B\setminus A_i)$.
  \item \textbf{Intersecci\'on de uniones:}
        $$\Big(\bigcup_{i\in I}A_i\Big)\cap\Big(\bigcup_{j\in J}B_j\Big)=\bigcup_{(i,j)\in I\times J}(A_i\cap B_j)$$
  \item \textbf{Uni\'on de intersecciones:}
        $$\Big(\bigcap_{i\in I}A_i\Big)\cup\Big(\bigcap_{j\in J}B_j\Big)=\bigcap_{(i,j)\in I\times J}(A_i\cup B_j)$$
  \item \textbf{No conmutan:}
        $\bigcup_i\bigcap_jA_{ij}\subseteq\bigcap_j\bigcup_iA_{ij}$, y la inclusi\'on inversa es falsa.
  \item \textbf{Disjuntificaci\'on:} con $B_i=A_i\setminus\bigcup_{j<i}A_j$ los $B_i$ son disjuntos dos a dos,
        $B_i\subseteq A_i$ y $\bigcup_iA_i=\bigsqcup_iB_i$; luego $\card{\bigcup A_i}=\sum\card{B_i}$.
  \item \textbf{Partici\'on:} si $\{P_k\}$ particiona $U$, $X=\bigsqcup_k(X\cap P_k)$ y $\card X=\sum_k\card{X\cap P_k}$.
  \item \textbf{Subaditividad:} $\card{\bigcup_iA_i}\le\sum_i\card{A_i}$.  Igualdad ssi disjuntos dos a dos.
  \item $\limsup A_n=\bigcap_n\bigcup_{k\ge n}A_k=\{x: x\in A_k\text{ para infinitos }k\}$.
  \item $\liminf A_n=\bigcup_n\bigcap_{k\ge n}A_k=\{x: x\in A_k\ \forall k\text{ salvo finitos}\}$
  \item $\liminf A_n\subseteq\limsup A_n$.
\end{itemize}
 
% =====================================================
\subsection{Inclusi\'on--exclusi\'on}
\begin{itemize}
  \item $\Big|\bigcup_{i=1}^nA_i\Big|=\sum_{\varnothing\ne S\subseteq[n]}(-1)^{|S|+1}\Big|\bigcap_{i\in S}A_i\Big|$.
  \item $\Big|U\setminus\bigcup_iA_i\Big|=\sum_{S\subseteq[n]}(-1)^{|S|}\Big|\bigcap_{i\in S}A_i\Big|$, con $\bigcap_\varnothing=U$.
  \item Sobreyecciones $[n]\to[m]$: $\sum_{k=0}^m(-1)^k\binom mk(m-k)^n$.
  \item Desarreglos: $D_n=\sum_{k=0}^n(-1)^k\,n!/k!$.
  \item Coprimos con $m$ hasta $N$: $\sum_{d\mid\rad(m)}\mu(d)\lfloor N/d\rfloor$. Donde $$m=p_1^{e_1}\cdots p_r^{e_r}\ \Longrightarrow\ \operatorname{rad}(m)=p_1p_2\cdots p_r.$$
\end{itemize}


% =====================================================
\subsection{Im\'agenes y preim\'agenes ($f:X\to Y$)}
\begin{itemize}
  \item $f^{-1}\big(\bigcup B_i\big)=\bigcup f^{-1}(B_i)$, \; $f^{-1}\big(\bigcap B_i\big)=\bigcap f^{-1}(B_i)$.
  \item $f^{-1}(\ov B)=\ov{f^{-1}(B)}$, \quad $f^{-1}(B\setminus C)=f^{-1}(B)\setminus f^{-1}(C)$.
  \item $f\big(\bigcup A_i\big)=\bigcup f(A_i)$.
  \item $f\big(\bigcap A_i\big)\subseteq\bigcap f(A_i)$.
  \item $f(A\setminus B)\supseteq f(A)\setminus f(B)$, con igualdad si $f$ es inyectiva.
  \item $f(f^{-1}(B))=B\cap f(X)$, \quad $f^{-1}(f(A))\supseteq A$.
\end{itemize}
 


 % =====================================================
\subsection{Producto cartesiano y conjunto potencia}
\begin{itemize}
  \item $(A\times B)\cap(C\times D)=(A\cap C)\times(B\cap D)$; \quad $\card{A\times B}=\card A\,\card B$.
  \item $A\times(B\cup C)=(A\times B)\cup(A\times C)$, \quad $A\times(B\cap C)=(A\times B)\cap(A\times C)$.
  \item $(A\times B)\cup(C\times D)\subsetneq(A\cup C)\times(B\cup D)$ en general.
  \item $\Pw(A)\cap\Pw(B)=\Pw(A\cap B)$.
  \item $\Pw(A)\cup\Pw(B)\subseteq\Pw(A\cup B)$, estricto salvo que $A\subseteq B$ o $B\subseteq A$.
\end{itemize}
 
% =====================================================
\subsection{Trucos frecuentes en problemas}
\begin{itemize}
  \item Elementos que aparecen una cantidad impar de veces: XOR (o $\sd$) de todos los $A_i$.
  \item Doble conteo: $\sum_i\card{A_i}=\sum_x\#\{i:x\in A_i\}$.
  \item Pares intersecantes: $\#\{(i,j):A_i\cap A_j\ne\varnothing\}=\text{total}-\#\{\text{disjuntos}\}$.
  \item Uni\'on de intervalos: ordenar y barrer; la longitud de la uni\'on \textbf{no} es la suma de longitudes (disjuntificar).
  \item Antes de intercambiar $\bigcup\bigcap$, revisa que los cuantificadores no conmutan (fuente t\'ipica de bugs).
  \item $\ov A$ y $\bigcap_\varnothing$ dependen de $U$; en m\'ascaras, el complemento exige \texttt{\& full}.
\end{itemize}



% =====================================================
\section{Subconjuntos y funciones sobre m\'ascaras}
\begin{itemize}
  \item $\card{\Pw(A)}=2^{\card A}$; hay $\binom nk$ subconjuntos de tama\~no $k$; $\sum_k\binom nk=2^n$; $\sum_k(-1)^k\binom nk=0$ para $n\ge1$.
  \item Pares $A\subseteq B\subseteq[n]$: $3^n$ (cada elemento: fuera de $B$, en $B\setminus A$, o en $A$).
        De aqu\'i el costo $O(3^n)$ de enumerar los submasks de todas las m\'ascaras.
  \item Cadenas $A_1\subseteq A_2\subseteq\dots\subseteq A_k\subseteq[n]$: $(k+1)^n$.
  \item \textbf{Zeta} (subconjuntos): $\hat f(S)=\sum_{T\subseteq S}f(T)$, en $O(n2^n)$ procesando un bit a la vez.
  \item Convoluci\'on OR: $\widehat{f*_{\cup}g}=\hat f\cdot\hat g$ (producto puntual tras la zeta).
  \item Convoluci\'on AND: igual, con la transformada sobre supersets. Convoluci\'on XOR: Walsh--Hadamard.
\end{itemize}
\newcommand{\piso}[1]{\lfloor#1\rfloor}

\subsection{Mobius General}
 Para cualquier conjunto con orden parcial $\leq$ existe una funcion $\mu(x,y)$ que depende solo del orden
\begin{center}
\small
\renewcommand{\arraystretch}{1.3}
\begin{tabular}{ll}
\hline
Suma hacia abajo & Suma hacia arriba \\
\hline \hline 
  \textit{Divisibilidad} \\
  $g(n)=\displaystyle\sum_{d\mid n}^{\piso{V/d}}f(d)$
  & $g(d)=\displaystyle\sum_{k\ge1}^{\piso{V/d}}f(kd)$ \\
  
  $f(n)=\displaystyle\sum_{d\mid n}\mu(d)\,g(n/d)$
  & $f(d)=\displaystyle\sum_{k\ge1}\mu(k)\,g(kd)$ \\
  \hline 
  \textit{Subconjuntos} \\ 
  $\hat f(S)=\displaystyle\sum_{T\subseteq S}f(T)$
  & $\check f(S)=\displaystyle\sum_{T\supseteq S}f(T)$ \\
  $f(S)=\displaystyle\sum_{T\subseteq S}(-1)^{|S\setminus T|}\hat f(T)$
  & $f(S)=\displaystyle\sum_{T\supseteq S}(-1)^{|T\setminus S|}\check f(T)$ \\
  \hline 
  \textit{General (POSET)}  \\ 
  $g(x)=\displaystyle\sum_{y\le x}f(y)$
  & $g(x)=\displaystyle\sum_{y\ge x}f(y)$ \\
  $f(x)=\displaystyle\sum_{y\le x}\mu(y,x)\,g(y)$
  & $f(x)=\displaystyle\sum_{y\ge x}\mu(x,y)\,g(y)$ \\
\hline
\end{tabular}
\end{center}
%=%=%=%=%=%=%=%=%=%=%=%=%=%=%=%=%=%=%=%=%=%=%=%=%=%=
%=%=%=%=%=%=%=%=%=%=%=%=%=%=%=%=%=%=%=%=%=%=%=%=%=%=
%=%=%=%=%=%=%=%=%=%=%=%=%=%=%=%=%=%=%=%=%=%=%=%=%=%=
%=%=%=%=%=%=%=%=%=%=%=%=%=%=%=%=%=%=%=%=%=%=%=%=%=%=
% \columnbreak
\section{Probabilidad y Procesos Estocásticos}


\subsection{Distribuciones Discretas Fundamentales}

\subsubsection{Distribución de Bernoulli}
Modela un único ensayo aleatorio con dos posibles resultados: Éxito ($1$) con probabilidad $p$, o Fracaso ($0$) con probabilidad $1-p$.
\begin{itemize}
    \item \textbf{Función de Masa de Probabilidad (PMF):} $P(X = k) = p^k (1-p)^{1-k}, \quad k \in \{0, 1\}$
    \item \textbf{Esperanza:} $E[X] = 1 \cdot p + 0 \cdot (1-p) = p$
    \item \textbf{Varianza:} $\text{Var}(X) = p(1-p)$
    \item \textbf{Variable Indicadora:} Si $X_i$ es la variable indicadora de que ocurra un evento $A$, entonces $E[X_i] = P(A)$.
\end{itemize}

\subsubsection{Distribución Binomial}
Modela la cantidad de éxitos $k$ obtenidos al realizar $n$ ensayos independientes de Bernoulli de probabilidad $p$.
\begin{itemize}
    \item \textbf{PMF:} $P(X = k) = \binom{n}{k} p^k (1-p)^{n-k}, \quad k \in \{0, 1, \dots, n\}$
    \item \textbf{Esperanza:} Por linealidad de la esperanza ($X = \sum_{j=1}^n X_j$):
    $$E[X] = \sum_{j=1}^n E[X_j] = n \cdot p$$
\end{itemize}

\subsubsection{Distribución Geométrica y Tiempo de Parada}
Modela el tiempo de parada $T$, definido como el número de ensayos de Bernoulli independientes con probabilidad de éxito $p$ necesarios hasta obtener el \textbf{primer éxito}.
\begin{itemize}
    \item \textbf{PMF:} $P(T = i) = (1-p)^{i-1} p, \quad i \in \{1, 2, 3, \dots\}$
    \item \textbf{Esperanza (Tiempo Esperado de Parada):}
    $$E[T] = \sum_{i=1}^{\infty} i \cdot P(T = i) = \sum_{i=1}^{\infty} i (1-p)^{i-1} p = \frac{1}{p}$$
    \item \textbf{Interpretación de Casos Favorables entre Casos Totales:}
    Si el espacio muestral en $i$ pasos tiene $M^i$ secuencias equiprobables y en los primeros $i-1$ pasos hay $F$ opciones de fracaso y en el $i$-ésimo hay $S$ opciones de éxito (donde $p = S/M$ y $1-p = F/M$):
    $$P(T = i) = \frac{F^{i-1} \cdot S}{M^i} = \left(\frac{F}{M}\right)^{i-1} \left(\frac{S}{M}\right)$$
\end{itemize}



\subsection{Técnicas de Análisis en Procesos Estocásticos y Cadenas de Markov}

\subsubsection{Análisis del Primer Paso (First-Step Analysis)}
Para encontrar el tiempo esperado de llegada $E[u]$ desde un estado $u$ hacia un estado terminal, se condiciona sobre la primera transición saliente:
$$E[u] = 1 + \sum_{v} P(u \to v) \cdot E[v]$$
Donde el $+1$ representa el costo/tiempo unitario consumido al realizar el paso actual.

\subsubsection{Despeje de Auto-ciclos (Self-Loops)}
Si un estado $u$ transiciona a sí mismo con probabilidad $p_{self} = P(u \to u)$:
$$E[u] = 1 + p_{self} \cdot E[u] + \sum_{v \neq u} P(u \to v) E[v]$$
Despejando $E[u]$ algebraicamente para evitar recursión infinita:
$$E[u] = \frac{1 + \sum_{v \neq u} P(u \to v) E[v]}{1 - p_{self}}$$

\subsubsection{Linealidad de la Esperanza y Descomposición}
Para cualquier conjunto de variables aleatorias $X_1, X_2, \dots, X_n$ (sean independientes o dependientes):
$$E\left[\sum_{j=1}^n X_j\right] = \sum_{j=1}^n E[X_j]$$
Permite calcular el tiempo esperado total de transformar $N$ elementos como $N$ veces el tiempo esperado de transformar un solo elemento, siempre que el proceso afecte a cada elemento de manera independiente.
%=%=%=%=%=%=%=%=%=%=%=%=%=%=%=%=%=%=%=%=%=%=%=%=%=%=
%=%=%=%=%=%=%=%=%=%=%=%=%=%=%=%=%=%=%=%=%=%=%=%=%=%=

\subsection{Ideas:}
\begin{description}
    \item[Maximo] $max_i \le c$ $\iff$  $\forall i \ X_i \le c$ 
    \item[] El proceso primero elige una situación o estado y luego evoluciona de manera distinta según esa elección.
    $$ \mathbb E[X]=\sum_s\Pr(S=s)\mathbb E[X\mid S=s]$$
\end{description}





%=%=%=%=%=%=%=%=%=%=%=%=%=%=%=%=%=%=%=%=%=%=%=%=%=%=
%=%=%=%=%=%=%=%=%=%=%=%=%=%=%=%=%=%=%=%=%=%=%=%=%=%=
%=%=%=%=%=%=%=%=%=%=%=%=%=%=%=%=%=%=%=%=%=%=%=%=%=%=
\columnbreak
\section{Plantilla Python}
\begin{minted}{cpp}
# librerias
import os, random, re, sys
from collections import deque
# limite memoria
sys.setrecursionlimit(200000)
# --- --- --- INTPUT & FAST --- --- ---
sys.stdin = open("in.txt", 'r', encoding='utf-8') # comentar
input = sys.stdin.readline
# --- --- --- DEBUG --- --- ---
DEB = False # True False
def debug(*args, **kwargs):
    if DEB:
        print("[DEBUG]", *args, **kwargs, file=sys.stderr)
# --- --- --- Leer y resolver --- --- ---
def solve():
    # separar multiples enteros en una línea
    n,m  = map(int, input().split())  
    # leer y quitar salto '\n'
    n = int( input().rstrip() ) 
    # leer una lista de enteros
    h = list( map(int, input().rstrip().split()) )
    # cadena de texto sin espacios
    s = input().rstrip()
    # escribir codigo aqui
# --- --- --- MAIN --- --- ---
if __name__ == "__main__":
    t = int( input() ) 
    for i in range(t):
        solve()   
\end{minted}


\subsection{Fenwick Tree}

Query puntual, actualizaciones en rangos. $O(\log n)$

\begin{minted}{cpp}   
class Fenwick:
    def __init__( self, n ):
        self.bit = [0] * (n+1)
        self.n = n
    def add( self, i, val ): # agrega puntual
        i += 1
        while i <= self.n:
            self.bit[i] += val
            i += i & (-i)  
    def add_rango( self, l, r, val): # en rango
        self.add( l, val)
        if r+ 1 < self.n :
            self.add( r+1, -val )
    def query_puntual( self, i): # query para un solo punto
        res = 0
        i += 1
        while i > 0:
            res += self.bit[i]
            i -= i & (-i)
        return res 
def solve():
    n,m  =   map(int, input().split())
    arbol = Fenwick(n)
    for _ in range(m):
        a,b, val =   map(int, input().split())
        arbol.add_rango( l=a-1, r=b-1, val = val)
    for i in range(n):
        print( arbol.query_puntual(i) )
\end{minted}


%%%%%%%%%%%%%%%%%%%%%%%%%%%%%%%%%%%%%%%%%%%%%%%%%%%%%%%%
%%%%%%%%%%%%%%%%%%%%%%%%%%%%%%%%%%%%%%%%%%%%%%%%%%%%%%%%
%%%%%%%%%%%%%%%%%%%%    PROBA    %%%%%%%%%%%%%%%%%%%%%%%
%%%%%%%%%%%%%%%%%%%%%%%%%%%%%%%%%%%%%%%%%%%%%%%%%%%%%%%%
%%%%%%%%%%%%%%%%%%%%%%%%%%%%%%%%%%%%%%%%%%%%%%%%%%%%%%%%
%%%%%%%%%%%%%%%%%%%%%%%%%%%%%%%%%%%%%%%%%%%%%%%%%%%%%%%%
\clearpage
\section{Modelado}

%%%%%%%%%%%%%%%%%%%%%%%%%%%%%%%%%%%%%%%%%%%%%%%%%%%%%%%%
%%                           PARTE 1
%%%%%%%%%%%%%%%%%%%%%%%%%%%%%%%%%%%%%%%%%%%%%%%%%%%%%%%%

\subsection{Distribuciones}
Discretas
\begin{description}
\item[Bernoulli] $X\!\sim\!\mathrm{Ber}(p)$, $p\!\in\!(0,1)$; $\mathcal S_X\!=\!\{0,1\}$; $f(x)\!=\!p^x(1\!-\!p)^{1-x}$; $\E\!=\!p$; $\Var\!=\!p(1\!-\!p)$; $M(t)\!=\!1\!-\!p\!+\!pe^t$.
\item[Binomial] $X\!\sim\!\mathrm{Bin}(n,p)$; $\mathcal S_X\!=\!\{0,\dots,n\}$; $f(x)\!=\!\binom nx p^x(1\!-\!p)^{n-x}$; $\E\!=\!np$; $\Var\!=\!np(1\!-\!p)$; $M(t)\!=\!(1\!-\!p\!+\!pe^t)^n$.
\item[Geométrica] (ensayos hasta el 1.er éxito) $X\!\sim\!\mathrm{Geo}(p)$, $p\!\in\!(0,1]$; $\mathcal S_X\!=\!\{1,2,\dots\}$; $f(x)\!=\!(1\!-\!p)^{x-1}p$; $\E\!=\!\frac1p$; $\Var\!=\!\frac{1-p}{p^2}$; $M(t)\!=\!\frac{pe^t}{1-(1-p)e^t}$, $t\!<\!-\ln(1\!-\!p)$.
\item[Bin. Negativa] (ensayos hasta el $r$-ésimo éxito) $X\!\sim\!\mathrm{NB}(r,p)$; $\mathcal S_X\!=\!\{r,r\!+\!1,\dots\}$; $f(x)\!=\!\binom{x-1}{r-1}p^r(1\!-\!p)^{x-r}$; $\E\!=\!\frac rp$; $\Var\!=\!\frac{r(1-p)}{p^2}$; $M(t)\!=\!\big(\frac{pe^t}{1-(1-p)e^t}\big)^r$, $t\!<\!-\ln(1\!-\!p)$.
\item[Poisson] $X\!\sim\!\mathrm{Poi}(\lambda)$, $\lambda\!>\!0$; $\mathcal S_X\!=\!\{0,1,\dots\}$; $f(x)\!=\!\frac{e^{-\lambda}\lambda^x}{x!}$; $\E\!=\!\Var\!=\!\lambda$; $M(t)\!=\!e^{\lambda(e^t-1)}$.
\item[Hipergeom.] $X\!\sim\!\mathrm{Hyp}(N,K,n)$; $\mathcal S_X\!=\!\{\max(0,n\!-\!N\!+\!K),\dots,\min(n,K)\}$; $f(x)\!=\!\frac{\binom Kx\binom{N-K}{n-x}}{\binom Nn}$; $\E\!=\!\frac{nK}N$; $\Var\!=\!n\frac KN(1\!-\!\frac KN)\frac{N-n}{N-1}$; $M(t)$ sin forma elemental (${}_2F_1$).
\end{description}
Continuas
\begin{description}
\item[Uniforme] $X\!\sim\!\mathcal U(a,b)$; $\mathcal S_X\!=\![a,b]$; $f(x)\!=\!\frac1{b-a}$; $\E\!=\!\frac{a+b}2$; $\Var\!=\!\frac{(b-a)^2}{12}$; $M(t)\!=\!\frac{e^{tb}-e^{ta}}{t(b-a)}$ ($t\!\neq\!0$), $M(0)\!=\!1$.
\item[Exponencial] $X\!\sim\!\mathrm{Exp}(\lambda)$, $\lambda\!>\!0$; $\mathcal S_X\!=\![0,\infty)$; $f(x)\!=\!\lambda e^{-\lambda x}$; $\E\!=\!\frac1\lambda$; $\Var\!=\!\frac1{\lambda^2}$; $M(t)\!=\!\frac\lambda{\lambda-t}$, $t\!<\!\lambda$.
\item[Normal] $X\!\sim\!\mathcal N(\mu,\sigma^2)$, $\sigma\!>\!0$; $\mathcal S_X\!=\!\mathbb R$; $f(x)\!=\!\frac1{\sigma\sqrt{2\pi}}e^{-(x-\mu)^2/2\sigma^2}$; $\E\!=\!\mu$; $\Var\!=\!\sigma^2$; $M(t)\!=\!e^{\mu t+\sigma^2t^2/2}$.
\item[Gamma] (forma--tasa) $X\!\sim\!\mathrm{Gamma}(\alpha,\beta)$; $\mathcal S_X\!=\!(0,\infty)$; $f(x)\!=\!\frac{\beta^\alpha}{\Gamma(\alpha)}x^{\alpha-1}e^{-\beta x}$; $\E\!=\!\frac\alpha\beta$; $\Var\!=\!\frac\alpha{\beta^2}$; $M(t)\!=\!(\frac\beta{\beta-t})^\alpha$, $t\!<\!\beta$.
\item[Beta] $X\!\sim\!\mathrm{Beta}(\alpha,\beta)$; $\mathcal S_X\!=\![0,1]$; $f(x)\!=\!\frac{x^{\alpha-1}(1-x)^{\beta-1}}{B(\alpha,\beta)}$, $B\!=\!\frac{\Gamma(\alpha)\Gamma(\beta)}{\Gamma(\alpha+\beta)}$; $\E\!=\!\frac\alpha{\alpha+\beta}$; $\Var\!=\!\frac{\alpha\beta}{(\alpha+\beta)^2(\alpha+\beta+1)}$; $M(t)\!=\!{}_1F_1(\alpha;\alpha\!+\!\beta;t)$.
\end{description}


%%%%%%%%%%%%%%%%%%%%%%%%%%%%%%%%%%%%%%%%%%%%%%%%%%%%%%%%
%%                           PARTE 2
%%%%%%%%%%%%%%%%%%%%%%%%%%%%%%%%%%%%%%%%%%%%%%%%%%%%%%%%


\subsection{Casos de uso y modelación}
\begin{description}
\item[Guía rápida: señal $\to$ modelo] ``$n$ intentos iguales e independientes, contar éxitos'' $\to$ Bin; ``intentos hasta el primer/$r$-ésimo éxito'' $\to$ Geo/NB; ``eventos raros en un intervalo'', sin $n$ explícito $\to$ Poi; ``sacar sin reemplazo'' $\to$ Hyp; ``punto/instante al azar'' $\to$ $\mathcal U$; ``tiempo hasta que ocurre'' con tasa constante $\to$ Exp; ``tiempo hasta el $k$-ésimo evento'' $\to$ Gamma; ``suma de muchas partes pequeñas'' $\to$ $\mathcal N$ (aproximación); ``$p$ desconocida'' o ``$p$ uniforme'' $\to$ Beta (Laplace); ``solo importa el orden relativo'' $\to$ permutación uniforme (intercambiabilidad). \textbf{Antes de modelar:} (i) ¿hay independencia o solo intercambiabilidad? (ii) ¿la linealidad de $\E$ resuelve sin conocer la distribución? (iii) ¿se pide exactitud o tolerancia $10^{-6}$?

\item[Ber/Bin] \textbf{Señal:} $n$ intentos independientes con la misma $p$; indicador de un evento. \textbf{Modelo:} $X=\sum_{j\le n}X_j$, $X_j\overset{iid}\sim\mathrm{Ber}(p)$; $\Prob(X\ge1)=1-(1-p)^n$ (complemento). \textbf{Relaciones:} $\mathrm{Bin}(n,p)+\mathrm{Bin}(m,p)=\mathrm{Bin}(n + m,p)$ (indep.); $\approx\mathrm{Poi}(np)$ si $n\to\infty$, $np\to\lambda$; $\approx\mathcal N(np,np(1 - p))$ si $np(1 - p)\gtrsim10$ (con corrección de continuidad $\pm\frac12$); si $X\sim\mathrm{Bin}(n,p)$, $Y\sim\mathrm{Bin}(m,p)$ indep., entonces $X\mid(X + Y{=}s)\sim\mathrm{Hyp}$. \textbf{En CP:} $\E[\sum X_j]=\sum\E[X_j]$ vale \emph{aunque haya dependencia} (la varianza sí exige independencia); pmf con recurrencia $f(k + 1)=f(k)\frac{n-k}{k+1}\frac p{1-p}$ en $O(n)$, o factoriales precalculados (mod $M$) en $O(1)$ por término; $n$ grande en reales $\Rightarrow$ logaritmos de factoriales (\texttt{lgamma}).

\item[Geo] \textbf{Señal:} ``cuántos intentos hasta el primer éxito'', ``se repite hasta que sale'', reinicio tras fallo. \textbf{Modelo:} tiempo de espera discreto; $\Prob(X > k)=(1 - p)^k$; es la única discreta sin memoria: $\Prob(X > m + k\mid X>m)=\Prob(X > k)$. \textbf{Relaciones:} $\min_i\mathrm{Geo}(p_i)\sim\mathrm{Geo}\bigl(1-\prod_i(1 - p_i)\bigr)$ (indep.); $\mathrm{Geo}(\lambda/n)/n\to\mathrm{Exp}(\lambda)$. \textbf{En CP:} en un DP, el autolazo se despeja como $E=c+(1 - p)E+\dots\Rightarrow E=\frac{c+\dots}{p}$ (no hay que iterar); muestreo por inversión $\lceil\ln U/\ln(1 - p)\rceil$.

\item[NB] \textbf{Señal:} ``intentos hasta acumular $r$ éxitos''. \textbf{Modelo:} $X=\sum_{j\le r}G_j$, $G_j\overset{iid}\sim\mathrm{Geo}(p)$; en la pmf el último ensayo es éxito y los $r - 1$ anteriores se reparten entre las $x - 1$ posiciones. \textbf{Relaciones:} dualidad $\Prob(\mathrm{NB}(r,p)\le x)=\Prob(\mathrm{Bin}(x,p)\ge r)$; versión por fracasos $=X-r$ (soporte desde 0; cuidar la convención). \textbf{En CP:} la dualidad evita un DP sobre ``éxitos acumulados'' y da la cola en $O(r)$ sumando una binomial.

\item[Poisson] \textbf{Señal:} conteo de eventos raros e independientes en un intervalo, dado el promedio $\lambda$ pero no $n$; llegadas a una cola. \textbf{Modelo:} $N(t)\sim\mathrm{Poi}(\lambda t)$ con conteos independientes en intervalos disjuntos. \textbf{Relaciones:} superposición: $\mathrm{Poi}(\lambda_1) + \mathrm{Poi}(\lambda_2)\sim\mathrm{Poi}(\lambda_1 + \lambda_2)$; \emph{thinning}: conservar cada evento con prob.\ $q$ da $\mathrm{Poi}(\lambda q)$, y conservados y descartados son independientes; $N_1\mid(N_1 + N_2{=}s)\sim\mathrm{Bin}\bigl(s,\frac{\lambda_1}{\lambda_1+\lambda_2}\bigr)$; $\mathrm{Bin}(n,p)\approx\mathrm{Poi}(np)$. \textbf{En CP:} recurrencia $f(k + 1)=f(k)\frac\lambda{k+1}$ y truncar la serie hacia $k\approx\lambda+c\sqrt\lambda$; \emph{Poissonización}: con $n$ fijo los conteos por categoría son dependientes, pero con $N\sim\mathrm{Poi}(\lambda)$ son independientes $\Rightarrow$ se calcula fácil y se des-poissoniza.

\item[Hyp] \textbf{Señal:} elegir $n$ de $N$ elementos distinguibles sin reemplazo (cartas, lotes, equipos). \textbf{Modelo:} cada elemento queda elegido con prob.\ $\frac nN$ por simetría, así que $\E[\#\text{rojos elegidos}]=\frac{nK}N$ sin importar la dependencia; la $n$-ésima extracción es roja con prob.\ $\frac KN$ (intercambiabilidad), pero \emph{condicionada a $g$ rojas previas} es $\frac{K-g}{N-(n-1)}$. \textbf{Relaciones:} si $N\to\infty$ y $K/N\to p$, $\to\mathrm{Bin}(n,p)$; versión multivariada para varias categorías. \textbf{En CP:} $\Prob=\binom Kk\binom{N-K}{n-k}/\binom Nn$ con factoriales modulares en $O(1)$; para probabilidades de configuraciones, contar casos favorables $/$ totales evita DP.

\item[Uniforme] \textbf{Señal:} ``punto/instante al azar'', sin preferencia; ``permutación aleatoria''. \textbf{Modelo:} probabilidad $\propto$ longitud, área o volumen (problemas geométricos se reducen a áreas en el cuadrado unidad); si solo importa el orden de $n$ uniformes iid continuas, todas las $n!$ permutaciones son equiprobables. \textbf{Relaciones:} $U\sim\mathcal U(0,1)\Rightarrow F^{-1}(U)\sim F$ (transformada inversa); $-\ln U\sim\mathrm{Exp}(1)$; estadísticos de orden $\to$ Beta. \textbf{En CP:} ``probabilidad de que el elemento $j$ sea el mayor/récord'' se resuelve por simetría sin integrar; para áreas complicadas, integral en $O(1)$ o partición en regiones.

\item[Exp] \textbf{Señal:} ``tiempo hasta que ocurre'' con tasa constante; varios procesos compitiendo en paralelo; ``falta de memoria''. \textbf{Modelo:} cadena de Markov en tiempo continuo: en el estado $s$ se espera $\mathrm{Exp}(\Lambda_s)$, $\Lambda_s=\sum\lambda_{s\to s'}$, y se salta a $s'$ con prob.\ $\frac{\lambda_{s\to s'}}{\Lambda_s}$. \textbf{Relaciones:} $\min_iT_i\sim\mathrm{Exp}(\sum\lambda_i)$, y $\Prob(T_k{=}\min)=\frac{\lambda_k}{\sum\lambda_i}$ (indep.); el máximo de $n$ iid $\mathrm{Exp}(\lambda)$ tiene $\E=\frac{H_n}\lambda$; los tiempos entre llegadas de un proceso de Poisson son Exp (dualidad). \textbf{En CP:} se convierte en DP discreto sobre la cadena embebida: $E[s]=\frac1{\Lambda_s}+\sum_{s'}\frac{\lambda_{s\to s'}}{\Lambda_s}E[s']$, sin integrales ni reales continuos.

\item[Normal] \textbf{Señal:} suma de muchas contribuciones pequeñas e independientes (TLC); enunciado que tolera error $\sim10^{-6}$ con $n$ enorme. \textbf{Modelo:} $\sum X_j\approx\mathcal N(n\mu,n\sigma^2)$; estandarizar con $Z=\frac{X-\mu}\sigma$ y usar $\Phi$. \textbf{Relaciones:} cerrada bajo suma independiente: $\mathcal N(\mu_1,\sigma_1^2)+\mathcal N(\mu_2,\sigma_2^2)=\mathcal N(\mu_1 + \mu_2,\sigma_1^2 + \sigma_2^2)$; aproxima Bin y Poi; Berry--Esseen: error $O(1/\sqrt n)$, malo en colas. \textbf{En CP:} es una \emph{aproximación}; si piden exactitud o respuesta modular no se usa. $\Phi(x)=\frac12\texttt{erfc}(-x/\sqrt2)$ en C++; para estimar tamaños o cotas es suficiente.

\item[Gamma] \textbf{Señal:} ``tiempo hasta el $k$-ésimo evento'' en un proceso de Poisson. \textbf{Modelo:} Erlang: $T=\sum_{j\le k}\mathrm{Exp}(\lambda)$. \textbf{Relaciones:} dualidad $\Prob(T\le t)=\Prob(\mathrm{Poi}(\lambda t)\ge k)$ (la integral se vuelve una suma finita de $k$ términos); $\mathrm{Gamma}(\alpha,\beta)+\mathrm{Gamma}(\alpha',\beta)=\mathrm{Gamma}(\alpha + \alpha',\beta)$ (indep., misma tasa); si $X\sim\mathrm{Gamma}(\alpha)$, $Y\sim\mathrm{Gamma}(\beta)$ indep., $\frac X{X+Y}\sim\mathrm{Beta}(\alpha,\beta)$ e independiente de $X + Y$. \textbf{En CP:} cola exacta en $O(k)$; normalizar $n$ exponenciales iid da un punto uniforme en el simplejo (Dirichlet(1,\dots,1)).

\item[Beta] \textbf{Señal:} probabilidad desconocida $p$ (``uniforme'' $\Rightarrow$ Beta(1,1)); posición del $k$-ésimo punto de $n$ uniformes. \textbf{Modelo:} prior conjugado de Ber/Bin con $a$ éxitos y $b$ fracasos virtuales: $p\mid(g\text{ de }i)\sim\mathrm{Beta}(a + g,b + i - g)$ y $\Prob(X_{i+1}{=}1\mid g\text{ de }i)=\frac{a+g}{a+b+i}$ (Laplace si $a{=}b{=}1$: $\frac{g+1}{i+2}$). \textbf{Relaciones:} $U_{(k)}\sim\mathrm{Beta}(k,n - k + 1)$; generalización a $k$ categorías: Dirichlet; urna de Pólya $\equiv$ Beta--Binomial. \textbf{En CP:} DP sobre $(i,g)$ con transición $\frac{a+g}{a+b+i}$ (inversos modulares de $a + b + i$) en $O(N^2)$; si solo importa el total de éxitos con Beta(1,1), es uniforme en $\{0,\dots,n\}$ y no hace falta DP.

\item[Mapa de relaciones] Ber $\xrightarrow{\Sigma}$ Bin $\to$ Poi (raro) $/\ \mathcal N$ (TLC); Geo $\xrightarrow{\Sigma}$ NB; Exp $\xrightarrow{\Sigma}$ Gamma; Geo $\xrightarrow{n\to\infty}$ Exp; Poi $\leftrightarrow$ Exp (conteo $\leftrightarrow$ espera); Bin $\leftrightarrow$ NB, Gamma $\leftrightarrow$ Poi (dualidades de colas); Hyp $\xrightarrow{N\to\infty}$ Bin; Gamma/Gamma $\to$ Beta; Beta $\leftrightarrow$ Bin (conjugación) y $U_{(k)}\sim$ Beta; Exp(1) $=-\ln\mathcal U$.
\end{description}


%%%%%%%%%%%%%%%%%%%%%%%%%%%%%%%%%%%%%%%%%%%%%%%%%%%%%%%%
%%                           PARTE 3
%%%%%%%%%%%%%%%%%%%%%%%%%%%%%%%%%%%%%%%%%%%%%%%%%%%%%%%%
\clearpage
\section{Modelacion avanzada}
\begin{description}

\item[Laplace / Beta--Binomial] \textbf{(1) Contexto.} Un evento (servidor falla, aparece un monstruo, moneda cae cara) ocurre con probabilidad fija pero desconocida $p$; tras $i$ intentos con $g$ éxitos se pide la probabilidad del siguiente, o un DP sobre estados $(i,g)$. \textbf{(2) Modelo.} $\theta\sim\mathrm{Beta}(a,b)$, $X_j\mid\theta\overset{\mathrm{iid}}{\sim}\mathrm{Ber}(\theta)$, $S_i=\sum_{j\le i}X_j$ (los $X_j$ son intercambiables, no independientes). 
\textbf{(3) Resultado.} 
$\theta\mid S_i{=}g\sim\mathrm{Beta}(a + g,\,b + i - g)$; $\Prob(X_{i+1}{=}1\mid S_i{=}g)=\frac{a+g}{a+b+i}$; $\Prob(S_n{=}k)=\binom nk\frac{B(a+k,\,b+n-k)}{B(a,b)}$; una secuencia concreta con $k$ éxitos tiene probabilidad $\frac{B(a+k,\,b+n-k)}{B(a,b)}$. Con $a{=}b{=}1$: $\frac{g+1}{i+2}$ y $S_n$ uniforme en $\{0,\dots,n\}$. \textbf{(4) Idea.} Integral beta $\int_0^1\theta^x(1 - \theta)^y d\theta=\frac{x!\,y!}{(x+y+1)!}$ (por partes) y Bayes: posterior $\propto$ verosimilitud $\times$ prior; Beta es conjugada de Bernoulli. Alternativa sin integrales: argumento de los huecos con $i + 2$ uniformes. \textbf{(5) CP.}  Conviene cuando ``$p$ fija desconocida'' y se aprende de los resultados.

\item[Dirichlet--Multinomial] \textbf{(1) Contexto.} Igual que Laplace pero con $k$ resultados posibles (dado de $k$ caras con sesgo desconocido, $k$ tipos de enemigos). \textbf{(2) Modelo.} $\theta\sim\mathrm{Dir}(\alpha_1,\dots,\alpha_k)$, $X_j\mid\theta\overset{\mathrm{iid}}{\sim}\mathrm{Cat}(\theta)$, conteos $c_m$ tras $i$ observaciones, $\alpha_0=\sum_m\alpha_m$. \textbf{(3) Resultado.} $\theta\mid c\sim\mathrm{Dir}(\alpha+c)$; $\Prob(X_{i+1}{=}m\mid c)=\frac{\alpha_m+c_m}{\alpha_0+i}$ (Laplace: $\frac{c_m+1}{i+k}$); $\Prob(c)=\frac{n!}{\prod_m c_m!}\cdot\frac{\Gamma(\alpha_0)}{\Gamma(\alpha_0+n)}\prod_m\frac{\Gamma(\alpha_m+c_m)}{\Gamma(\alpha_m)}$. \textbf{(4) Idea.} Integral de Dirichlet $\int_{\Delta}\prod\theta_m^{x_m-1}d\theta=\frac{\prod\Gamma(x_m)}{\Gamma(\sum x_m)}$ y conjugación; equivale a una urna de Pólya con $\alpha_m$ bolas de cada color. \textbf{(5) CP.} $O(k)$ por consulta o $O(1)$ manteniendo $i$ y $c$; el estado de un DP es el vector de conteos (compresión: contar solo cuántos colores tienen cada conteo). Conviene con varias categorías y probabilidades desconocidas e intercambiables.

\item[Urna de P\'olya] \textbf{(1) Contexto.} Se saca una bola, se devuelve y se agregan $c$ bolas del mismo color (``el rico se hace m\'as rico''). \textbf{(2) Modelo.} Inicio con $a$ rojas y $b$ azules; $R_n$ = rojas obtenidas en $n$ extracciones. \textbf{(3) Resultado.} $\Prob(\text{roja en }n + 1\mid R_n{=}r)=\frac{a+cr}{a+b+cn}$; una secuencia con $r$ rojas y $n - r$ azules tiene probabilidad $\frac{a^{(r,c)}\,b^{(n-r,c)}}{(a+b)^{(n,c)}}$ con $x^{(m,c)}=x(x + c)\cdots(x + (m - 1)c)$. Es exactamente Beta--Binomial con $(a/c,\,b/c)$ y la fracción límite es $\mathrm{Beta}(a/c,b/c)$. \textbf{(4) Idea.} Inducción sobre la secuencia: el numerador y el denominador son productos crecientes que no dependen del orden (intercambiabilidad). \textbf{(5) CP.} $O(1)$ por paso, DP $(n,r)$ en $O(N^2)$. Conviene cuando la probabilidad de la siguiente extracci\'on es proporcional al conteo acumulado.

\item[Estad\'isticos de orden y espaciamientos] \textbf{(1) Contexto.} Se eligen $n$ puntos al azar en $[0,1]$ (o una vara se parte en $n + 1$ trozos) y se pregunta por el mayor, el menor, el $k$-\'esimo o el trozo m\'as largo. \textbf{(2) Modelo.} $U_1,\dots,U_n\overset{\mathrm{iid}}{\sim}\mathcal U(0,1)$, $U_{(1)}\le\dots\le U_{(n)}$, espaciamientos $D_j=U_{(j)}-U_{(j-1)}$ con $U_{(0)}{=}0$, $U_{(n+1)}{=}1$. \textbf{(3) Resultado.} $U_{(k)}\sim\mathrm{Beta}(k,\,n - k + 1)$; $\E[U_{(k)}]=\frac k{n+1}$; $\Var=\frac{k(n-k+1)}{(n+1)^2(n+2)}$; $\E[\max]=\frac n{n+1}$, $\E[\min]=\frac1{n+1}$, $\E[\text{rango}]=\frac{n-1}{n+1}$; $(D_1,\dots,D_{n+1})\sim\mathrm{Dir}(1,\dots,1)$ (intercambiables, cada una con esperanza $\frac1{n+1}$); $\E[\max_j D_j]=\frac1{n+1}\sum_{j=1}^{n+1}\frac1j$; $\Prob(U_{n+1}\text{ es el }k\text{-\'esimo menor})=\frac1{n+1}$. Para otra $F$ continua: aplicar $F^{-1}$. \textbf{(4) Idea.} $\Prob(U_{(k)}\in dx)=\frac{n!}{(k-1)!(n-k)!}x^{k-1}(1 - x)^{n-k}dx$ (elegir cu\'ales $k - 1$ caen antes de $x$, una cae en $dx$, el resto despu\'es); espaciamientos intercambiables por simetr\'ia, y $\sum D_j=1$ da la esperanza. \textbf{(5) CP.} $O(1)$. Conviene cuando aparecen m\'aximos/m\'inimos de uniformes, partir intervalos al azar o ``probabilidad de que el nuevo punto sea r\'ecord''.

\item[Suma de colas y m\'aximo discreto] \textbf{(1) Contexto.} Esperanza del m\'aximo de $n$ dados de $d$ caras, o de cualquier v.a. entera no negativa con cola conocida. \textbf{(2) Modelo.} $X\ge0$ entera (o continua). \textbf{(3) Resultado.} $\E[X]=\sum_{k\ge1}\Prob(X\ge k)$ y $\E[X]=\int_0^\infty\Prob(X>t)\,dt$; para $M=\max$ de $n$ dados iid de $d$ caras: $\E[M]=\sum_{x=1}^{d}\bigl(1-(\tfrac{x-1}d)^n\bigr)=d-\sum_{x=1}^{d-1}(\tfrac xd)^n$. \textbf{(4) Idea.} $X=\sum_{k\ge1}\mathbf 1[X\ge k]$ y linealidad (o Fubini en el caso continuo); $\Prob(M\le x)=(x/d)^n$ por independencia. \textbf{(5) CP.} $O(d\log n)$ con exponenciaci\'on r\'apida. Conviene siempre que $\Prob(X\ge k)$ es m\'as f\'acil que la pmf.

\item[Coleccionista de cupones generalizado] \textbf{(1) Contexto.} Hay $n$ tipos de cup\'on y cada compra da el tipo $i$ con probabilidad $p_i$; se pide el n\'umero esperado de compras hasta tener todos (o $m$ distintos). \textbf{(2) Modelo.} $T$ = primer instante en que todos los tipos aparecieron, $\sum p_i=1$. \textbf{(3) Resultado.} $\E[T]=\int_0^\infty\Bigl(1-\prod_{i=1}^n(1-e^{-p_it})\Bigr)dt=\sum_{\emptyset\ne J\subseteq[n]}\frac{(-1)^{|J|+1}}{p(J)}$, con $p(J)=\sum_{i\in J}p_i$. Caso uniforme: $\E[T]=nH_n$ y $\Var(T)=n^2\sum_{j=1}^n\frac1{j^2}-nH_n$. Para $m$ distintos (uniforme): $\E=\sum_{k=0}^{m-1}\frac n{n-k}$. \textbf{(4) Idea.} Poissonizaci\'on: las compras llegan como proceso de Poisson de tasa 1, el tipo $i$ aparece como proceso independiente de tasa $p_i$, as\'i que el tiempo continuo es $\max_i\mathrm{Exp}(p_i)$ con $\Prob(\le t)=\prod_i(1 - e^{-p_it})$; el tiempo medio entre llegadas es 1 y por Wald $\E[T_{\text{cont}}]=\E[T]$. Expandir el producto da inclusi\'on--exclusi\'on con $\int_0^\infty e^{-p(J)t}dt=\frac1{p(J)}$. \textbf{(5) CP.} Inclusi\'on--exclusi\'on $O(2^n)$ ($n\lesssim20$); si los $p_i$ son enteros peque\~nos, DP sobre la suma $p(J)$ con signo en $O(n\cdot S)$; uniforme $O(n)$. Conviene con probabilidades no uniformes y $n$ peque\~no, o con DP por estados ``cu\'antos faltan''.

\item[Identidad Min--Max] \textbf{(1) Contexto.} Esperanza del tiempo en que \emph{todos} los elementos de un conjunto han ocurrido (cupones, procesos paralelos, lámparas que se apagan). \textbf{(2) Modelo.} Variables aleatorias $X_i$ (cualquier dependencia), $S=\{1,\dots,n\}$. \textbf{(3) Resultado.} $\max S=\sum_{\emptyset\ne T\subseteq S}(-1)^{|T|+1}\min T$, y por linealidad $\E[\max S]=\sum_{T\ne\emptyset}(-1)^{|T|+1}\E[\min T]$. Generalizaci\'on: $\E[k\text{-\'esimo mayor}]=\sum_{T}(-1)^{|T|-k}\binom{|T|-1}{k-1}\E[\min T]$. \textbf{(4) Idea.} Para reales $x_i\ge0$ y $t\ge0$: $\mathbf 1[\max>t]=1-\prod_i(1-\mathbf 1[x_i>t])=\sum_{T\ne\emptyset}(-1)^{|T|+1}\prod_{i\in T}\mathbf 1[x_i>t]=\sum_T(-1)^{|T|+1}\mathbf 1[\min T>t]$ (producto = indicador del m\'inimo), e integrar en $t$ (o aplicar cualquier funcional lineal). \textbf{(5) CP.} $O(2^n)$ en general; \'util cuando $\E[\min T]$ tiene forma cerrada (p. ej.\ $\frac1{p(T)}$ para exponenciales o cupones, con DP sobre la suma como arriba). Conviene si el m\'inimo de un subconjunto es f\'acil pero el m\'aximo no.

\item[Carrera de exponenciales y Poissonizaci\'on] \textbf{(1) Contexto.} Varios procesos independientes compiten (cola de servidores, cupones con tiempos continuos); se pide qui\'en termina primero o cu\'anto tarda el \'ultimo. \textbf{(2) Modelo.} $T_i\sim\mathrm{Exp}(\lambda_i)$ independientes, $\Lambda=\sum\lambda_i$. \textbf{(3) Resultado.} $\min_iT_i\sim\mathrm{Exp}(\Lambda)$; $\Prob(T_k=\min_iT_i)=\frac{\lambda_k}{\Lambda}$; el m\'inimo es independiente de cu\'al lo alcanza; si $\lambda_i{=}\lambda$: $\E[\max]=\frac{H_n}{\lambda}$ y $\E[k\text{-\'esimo en terminar}]=\frac1\lambda\sum_{j=0}^{k-1}\frac1{n-j}$. En una CTMC, cada estado dura $\mathrm{Exp}(\Lambda_s)$ con $\E=\frac1{\Lambda_s}$ y salta a $s'$ con probabilidad $\frac{\lambda_{s\to s'}}{\Lambda_s}$. \textbf{(4) Idea.} $\Prob(\min>t)=\prod e^{-\lambda_it}=e^{-\Lambda t}$; falta de memoria $\Rightarrow$ tras cada llegada el resto sigue siendo exponencial con los mismos par\'ametros. \textbf{(5) CP.} $O(1)$ por paso del DP sobre estados; convierte problemas continuos en cadenas discretas ponderadas por $\frac1\Lambda$. Conviene cuando el enunciado habla de tiempos continuos o tasas.

\item[Ruina del jugador] \textbf{(1) Contexto.} Un caminante empieza en $i$, sube con probabilidad $p$ y baja con $q=1 - p$, y se detiene al llegar a $0$ o a $N$. \textbf{(2) Modelo.} Paseo $S_n$ con barreras absorbentes; $r=q/p$; $P_i=\Prob(\text{llegar a }N\text{ antes que a }0)$; $T$ = tiempo de parada. \textbf{(3) Resultado.} $P_i=\frac{1-r^i}{1-r^N}$ si $p\ne q$ y $P_i=\frac iN$ si $p=q$; $\E[T_i]=i(N - i)$ si $p=q$, y $\E[T_i]=\frac1{q-p}\Bigl(i-N\,\frac{1-r^i}{1-r^N}\Bigr)$ si $p\ne q$. Con una sola barrera (en $0$): ruina con probabilidad $\min(1,r^i)$. \textbf{(4) Idea.} An\'alisis del primer paso: $P_i=pP_{i+1}+qP_{i-1}$, ecuaci\'on caracter\'istica $px^2-x+q=0$ con ra\'ices $1$ y $r$; para el tiempo $E_i=1+pE_{i+1}+qE_{i-1}$ con soluci\'on particular $\frac i{q-p}$. Alternativa: $r^{S_n}$ y $S_n-(p-q)n$ son martingalas; teorema de parada opcional. \textbf{(5) CP.} $O(1)$ con potencias modulares (si $r\equiv1$ mod $M$ usar el caso $p=q$ con cuidado). Conviene en pasos aleatorios 1D con barreras; con pasos no unitarios pasa a ecuaciones lineales (siguiente \'item).

\item[Cadenas de Markov absorbentes] \textbf{(1) Contexto.} Un proceso salta entre estados con probabilidades fijas hasta caer en un estado terminal; se pide la probabilidad de terminar en cada uno o el tiempo esperado. \textbf{(2) Modelo.} $P=\begin{pmatrix}Q&R\\0&I\end{pmatrix}$ con $Q$ transitorio$\to$transitorio y $R$ transitorio$\to$absorbente. \textbf{(3) Resultado.} Matriz fundamental $N=(I-Q)^{-1}=\sum_{k\ge0}Q^k$ ($N_{uv}$ = visitas esperadas a $v$ desde $u$); tiempos esperados $t=N\mathbf 1$; probabilidades de absorci\'on $B=NR$; probabilidad de estar en $s$ tras exactamente $T$ pasos: $e_u^\top P^T$. \textbf{(4) Idea.} $Q^k\to0$ (de cada estado transitorio se escapa con probabilidad positiva), luego la serie geom\'etrica matricial converge y $(I - Q)N=I$; sumar sobre columnas da el n\'umero esperado de pasos. \textbf{(5) CP.} Eliminaci\'on gaussiana $O(s^3)$ (modular: verificar que $I - Q$ sea invertible mod $M$); matrices tridiagonales (1D) en $O(s)$; horizonte $T$ grande con exponenciaci\'on $O(s^3\log T)$ o Kitamasa/Berlekamp--Massey si hay recurrencia lineal; en \'arboles/DAG escribir $E_v=a_vE_{\mathrm{padre}}+b_v$ y despejar en $O(n)$ (``inc\'ognita lineal''). Conviene en cualquier ``esperanza de pasos hasta terminar'' con ciclos.

\item[Paseos aleatorios en grafos] \textbf{(1) Contexto.} Una part\'icula salta a un vecino uniforme; se pide tiempo esperado para llegar a $v$, para volver, o para visitar todo. \textbf{(2) Modelo.} Grafo conexo no dirigido con $m$ aristas; $h(u,v)=\E[\text{pasos }u\to v]$. \textbf{(3) Resultado.} Distribuci\'on estacionaria $\pi_v=\frac{\deg v}{2m}$; tiempo de retorno $\E_v[T_v^+]=\frac{2m}{\deg v}$; tiempo de conmutaci\'on $h(u,v)+h(v,u)=2m\,R_{\mathrm{ef}}(u,v)$; tiempo de cobertura $\le2m(n - 1)$; en un camino de $n$ nodos con reflexi\'on en $0$, $h(0\to n)=n^2$. \textbf{(4) Idea.} $\pi$ satisface el balance detallado $\pi_u/\deg u=\pi_v/\deg v$ y la esperanza del retorno es $1/\pi_v$ (teorema de renovaci\'on); la f\'ormula de conmutaci\'on sale de la analog\'ia con redes el\'ectricas (leyes de Kirchhoff). \textbf{(5) CP.} $O(n^3)$ resolviendo el sistema $h(u,v)=1+\frac1{\deg u}\sum_{w\sim u}h(w,v)$, o $O(n)$ en \'arboles. Conviene en problemas de ``caminata aleatoria sobre un grafo/\'arbol''.

\item[Reflexi\'on y teorema del escrutinio] \textbf{(1) Contexto.} Dos candidatos reciben $a>b$ votos en orden aleatorio; probabilidad de que el ganador vaya siempre adelante (o nunca atr\'as); conteo de caminos que no tocan una barrera. \textbf{(2) Modelo.} Todas las $\binom{a+b}{a}$ secuencias equiprobables. \textbf{(3) Resultado.} $\Prob(\text{A siempre estrictamente adelante})=\frac{a-b}{a+b}$; $\Prob(\text{A nunca atr\'as})=\frac{a-b+1}{a+1}$; en particular, para $a{=}b{=}n$ la probabilidad de no pasar bajo cero es $\frac1{n+1}$ (n\'umeros de Catalan). \textbf{(4) Idea.} Reflexi\'on: los caminos de $A$ a $B$ que tocan la recta $y=-1$ se biyectan con los caminos de $A'$ (reflejado) a $B$, y $\binom{a+b}{a}-\binom{a+b}{a+1}$ cuenta los buenos; alternativa por rotaci\'on c\'iclica (lema de Dvoretzky--Motzkin). \textbf{(5) CP.} $O(1)$ con factoriales precalculados. Conviene en caminos con restricci\'on de signo o de barrera, y para depurar un DP $O(n^2)$ que cuenta lo mismo.

\item[Linealidad de la esperanza en grafos aleatorios] \textbf{(1) Contexto.} Esperanza del n\'umero de componentes, aristas de un \'arbol aleatorio, o grados, cuando el grafo se subsamplea al azar. \textbf{(2) Modelo.} Indicadores $\mathbf 1_e$, $\mathbf 1_v$ sin necesidad de independencia entre ellos. \textbf{(3) Resultado.} Bosque con $n$ v\'ertices, aristas conservadas con prob. $p$: $\E[\#\text{comp}]=n-(n - 1)p$; v\'ertices conservados con prob. $p$: $\E[\#\text{comp}]=np-(n - 1)p^2$ (componentes $=V'-E'$ en un bosque). $G(n,p)$: $c(n)=1-\sum_{k=1}^{n-1}\binom{n-1}{k-1}c(k)(1 - p)^{k(n-k)}$ es $\Prob(\text{conexo})$ y $\E[\#\text{comp}]=\sum_k\binom nk c(k)(1 - p)^{k(n-k)}$. \'Arbol de expansi\'on uniforme: $\tau(G)=\det$ de cualquier cofactor del Laplaciano, y $\Prob(e\in T)=R_{\mathrm{ef}}(e)$. \'Arbol etiquetado uniforme ($n^{n-2}$, Cayley): $\deg v-1\sim\mathrm{Bin}(n - 2,\tfrac1n)$, $\E[\#\text{hojas}]=n(1 - \tfrac1n)^{n-2}\approx\tfrac ne$. \textbf{(4) Idea.} Para $c(n)$: fijar el componente del v\'ertice $1$ (tama\~no $k$, conexo, sin aristas hacia el resto); linealidad $\E[\sum\mathbf 1]=\sum\Prob$; Kirchhoff y propiedad de Prüfer. \textbf{(5) CP.} $O(n^2)$ para $c(n)$, $O(n^3)$ para el determinante, $O(n)$ para bosques. Conviene cuando el enunciado pide ``esperanza de cu\'antos...'' (siempre intentar descomponer en indicadores antes de un DP).

\item[Linealidad en permutaciones aleatorias] \textbf{(1) Contexto.} Permutaci\'on uniforme de $n$ elementos; se piden puntos fijos, ciclos, r\'ecords o inversiones esperadas. \textbf{(2) Modelo.} $\sigma$ uniforme en $S_n$. \textbf{(3) Resultado.} $\E[\#\text{puntos fijos}]=1$; $\E[\#\text{ciclos}]=H_n$; $\E[\#\text{r\'ecords}]=H_n$ (el $i$-\'esimo es r\'ecord con prob. $\frac1i$, independientes); $\E[\#\text{inversiones}]=\frac{n(n-1)}4$; $\E[\text{longitud del ciclo de }1]=\frac{n+1}2$; $\Prob(\text{sin puntos fijos})=\frac{D_n}{n!}\to\frac1e$ con $D_n=n!\sum_{k=0}^n\frac{(-1)^k}{k!}$. \textbf{(4) Idea.} Indicadores $\mathbf 1[\sigma(i){=}i]$ con $\Prob=\frac1n$; $\mathbf 1[\text{par }(i,j)\text{ invertido}]$ con prob. $\frac12$ por simetr\'ia; la longitud del ciclo de $1$ es uniforme en $\{1,\dots,n\}$; inclusi\'on--exclusi\'on para $D_n$. \textbf{(5) CP.} $O(1)$ o $O(n)$ para sumar. Conviene siempre que haya simetr\'ia entre elementos.

\item[Rachas, patrones y f\'ormula de Conway] \textbf{(1) Contexto.} Esperanza de lanzamientos hasta ver $k$ \'exitos seguidos, o hasta ver un patr\'on $s$ (HHT, ABRACADABRA); probabilidad de que un patr\'on aparezca antes que otro. \textbf{(2) Modelo.} Ensayos Bernoulli($p$) (o letras iid), $T_s$ = primer tiempo en que $s$ aparece. \textbf{(3) Resultado.} Racha de $k$ \'exitos: $\E[T]=\sum_{j=1}^k p^{-j}=\frac{1-p^k}{(1-p)p^k}$. Patr\'on $s$ con letras iid: $\E[T_s]=\sum_{\ell:\ \mathrm{pref}_\ell(s)=\mathrm{suf}_\ell(s)}\frac1{\Prob(\mathrm{pref}_\ell(s))}$ (suma sobre los bordes de $s$, incluido $\ell{=}|s|$); moneda justa: $\sum 2^\ell$. Penney (moneda justa): con $L(X,Y)=\sum_k2^{k-1}[\text{los \'ultimos }k\text{ de }X=\text{los primeros }k\text{ de }Y]$, $\Prob(A\text{ antes que }B)=\frac{L(B,B)-L(B,A)}{(L(A,A)-L(A,B))+(L(B,B)-L(B,A))}$. \textbf{(4) Idea.} Racha: $E_k=E_{k-1}+1+(1 - p)E_k$ (si falla, se reinicia). Patr\'on: casino con apuestas justas, un jugador entra en cada instante y apuesta todo a la siguiente letra; el capital total al detener es $\sum_\ell1/\Prob(\mathrm{pref}_\ell)$ para los bordes (martingala y parada opcional). \textbf{(5) CP.} $O(k)$ o $O(1)$ con serie geom\'etrica; patrones: bordes con KMP/Z en $O(|s|)$; ``primero en aparecer'' con dos patrones se resuelve con $L$ en $O(|A|+|B|)$. Conviene en problemas con reinicio condicional (falla $\Rightarrow$ volver al inicio de un autómata).

\item[Reinicios geom\'etricos e identidad de Wald] \textbf{(1) Contexto.} Se repite un intento (que cuesta tiempo aleatorio) hasta que uno tiene \'exito; se pide el costo total esperado. \textbf{(2) Modelo.} Intentos iid con costo $C_i$ (incluye fallos) y \'exito con prob. $q>0$; $N$ = n\'umero de intentos, $T=\sum_{i\le N}C_i$. \textbf{(3) Resultado.} $\E[N]=\frac1q$ y $\E[T]=\E[N]\,\E[C]=\frac{\E[C]}{q}$; en general, para un tiempo de parada $\tau$ con $\E[\tau]<\infty$: $\E\bigl[\sum_{i\le\tau}X_i\bigr]=\E[\tau]\E[X_1]$. \textbf{(4) Idea.} $\sum_{i\le\tau}X_i=\sum_{i\ge1}X_i\mathbf 1[\tau\ge i]$ y $\{\tau\ge i\}$ es independiente de $X_i$ (depende solo de $X_1,\dots,X_{i-1}$); tomar esperanza. Contraejemplo: si $\tau$ ``mira hacia adelante'' (p. ej.\ $\tau$ = \'ultimo \'exito), Wald falla. \textbf{(5) CP.} $O(1)$. Conviene en ``estrategia con reintentos'' y para cerrar DP con autoloop: $E=c+(1 - q)E+\dots\Rightarrow E=\frac{c+\dots}{q}$.

\item[Colisiones y paradoja del cumplea\~nos] \textbf{(1) Contexto.} Cu\'antos elementos caben en $N$ casillas (hash de 64 bits, c\'odigos, cumplea\~nos) antes de que haya colisi\'on. \textbf{(2) Modelo.} $n$ valores uniformes iid en $\{1,\dots,N\}$. \textbf{(3) Resultado.} $\Prob(\text{sin colisi\'on})=\prod_{i=0}^{n-1}\bigl(1-\frac iN\bigr)\le e^{-n(n-1)/(2N)}$; $\E[\#\text{pares colisionados}]=\binom n2\frac1N$; $n\approx1.18\sqrt N$ da probabilidad $\tfrac12$; $\E[\text{tama\~no hasta 1.ª colisi\'on}]\approx\sqrt{\pi N/2}$. \textbf{(4) Idea.} Producto de probabilidades condicionales, $1-x\le e^{-x}$, y linealidad con indicadores de pares. \textbf{(5) CP.} $O(1)$ para estimar error de hashing (p. ej.\ $n=10^6$ con $N=2^{64}$: error $\approx\frac{n^2}{2N}\approx3\cdot10^{-8}$). Conviene al elegir m\'odulos/bases de hashing y al analizar algoritmos aleatorizados.

\item[Test polinomial (Schwartz--Zippel) y verificaci\'on aleatoria] \textbf{(1) Contexto.} Comprobar si dos expresiones/matrices son iguales (identidad polin\'omica, $AB=C$, igualdad de multiconjuntos) sin calcularlas por completo. \textbf{(2) Modelo.} $P\not\equiv0$ de grado total $d$ sobre un cuerpo $F$; $r$ uniforme en $S^k$, $S\subseteq F$. \textbf{(3) Resultado.} $\Prob(P(r)=0)\le\frac d{|S|}$; Freivalds: si $AB\ne C$ y $x\in\{0,1\}^n$ uniforme, $\Prob(ABx=Cx)\le\frac12$ (error $2^{-t}$ con $t$ repeticiones). \textbf{(4) Idea.} Inducci\'on en el n\'umero de variables: un polinomio univariado no nulo de grado $d$ tiene $\le d$ ra\'ices; fijar las dem\'as variables y condicionar. \textbf{(5) CP.} Freivalds $O(n^2)$ por prueba frente a $O(n^3)$; hashing polinomial de cadenas con base aleatoria. Conviene para verificar resultados caros con error unilateral controlado.

\item[Problema de la secretaria (parada \'optima)] \textbf{(1) Contexto.} Se ven $n$ candidatos en orden aleatorio y hay que aceptar o rechazar al instante; se quiere maximizar la probabilidad de elegir el mejor. \textbf{(2) Modelo.} Orden uniforme; estrategia: rechazar los primeros $r - 1$ y aceptar el primero mejor que todos los anteriores. \textbf{(3) Resultado.} $\Prob(\text{\'exito}_r)=\frac{r-1}n\sum_{j=r}^n\frac1{j-1}$ para $r\ge2$; el \'optimo es $r^\ast\approx\frac ne$ con probabilidad $\to\frac1e$. \textbf{(4) Idea.} El mejor est\'a en la posici\'on $j$ con prob. $\frac1n$, y se acepta si el mejor de los $j - 1$ anteriores est\'a entre los primeros $r - 1$ (prob. $\frac{r-1}{j-1}$); optimizar sobre $r$. \textbf{(5) CP.} $O(n)$ para probar todos los $r$; para variantes (valores conocidos, costos) se usa DP hacia atr\'as sobre la ``esperanza de seguir''. Conviene en parada \'optima: $V_i=\max(\text{aceptar},\E[V_{i+1}])$.

\item[Bayes con hip\'otesis discretas] \textbf{(1) Contexto.} Hay $k$ ``tipos'' posibles del sistema (dado justo/cargado, jugador fuerte/d\'ebil) y se observan resultados. \textbf{(2) Modelo.} Prior $\pi_h$, verosimilitudes $\ell_h(x)$. \textbf{(3) Resultado.} $\Prob(h\mid x_{1:i})=\frac{\pi_h\prod_{j\le i}\ell_h(x_j)}{\sum_{h'}\pi_{h'}\prod_{j\le i}\ell_{h'}(x_j)}$ y predicci\'on $\Prob(x_{i+1}\mid x_{1:i})=\sum_h\Prob(h\mid x_{1:i})\,\ell_h(x_{i+1})$. \textbf{(4) Idea.} Ley de probabilidad total sobre las hip\'otesis y regla de Bayes; el denominador es la verosimilitud marginal. \textbf{(5) CP.} $O(k)$ por observaci\'on; trabajar con logaritmos o renormalizar para evitar underflow, y para respuestas modulares multiplicar por inversos. Conviene cuando el enunciado dice ``con probabilidad $\pi_h$ el sistema es de tipo $h$'' (mezcla discreta, el an\'alogo discreto de Laplace).

\item[Otros patrones que conviene tener a mano]  \textbf{Esperanza condicionada:} $\E[X]=\E[\E[X\mid Y]]$ para dividir en casos; \textbf{varianza total:} $\Var X=\E[\Var(X\mid Y)]+\Var(\E[X\mid Y])$. \textbf{Funciones generadoras:} $G_X(z)=\E[z^X]$ convierte sumas independientes en productos (convoluci\'on con FFT/NTT) y da $\E[X]=G_X'(1)$. \textbf{Cotas de cola:} Markov $\Prob(X\ge a)\le\frac{\E X}a$, Chebyshev $\Prob(|X-\mu|\ge k\sigma)\le\frac1{k^2}$, Chernoff para sumas de Bernoulli: $\Prob(X\ge(1 + \delta)\mu)\le e^{-\delta^2\mu/3}$ ($0<\delta\le1$), para justificar algoritmos aleatorizados y heur\'isticas. \textbf{Lema de Borel--Cantelli / m\'etodo probabil\'istico:} si $\E[X]=\mu$ existe un objeto con $X\ge\mu$ y otro con $X\le\mu$ (cotas de tipo ``siempre existe un corte con $\ge m/2$ aristas''). \textbf{Desaleatorizaci\'on:} m\'etodo de esperanzas condicionales, escoger cada decisi\'on para no empeorar $\E[X\mid\text{decisiones previas}]$ y obtener un algoritmo determinista.

\end{description}
\clearpage



\end{multicols*}


\end{document}
%%%%%%%%%%%%%%%%%%%%%%%%%%%%%%%%%%%%%%%%%%%%%%%%%%%%%%%%
%%%%%%%%%%%%%%%%%%%%%%%%%%%%%%%%%%%%%%%%%%%%%%%%%%%%%%%%
%%%%%%%%%%%%%%%%%%%%%%%%%%%%%%%%%%%%%%%%%%%%%%%%%%%%%%%%
%%%%%%%%%%%%%%%%%%%%%%%%%%%%%%%%%%%%%%%%%%%%%%%%%%%%%%%%

\section*{Fundamentos del Espacio Probabilístico (Kolmogorov)}
Para cualquier problema, operamos sobre un espacio de probabilidad $(\Omega, \mathcal{F}, \mathbb{P})$, donde:
\begin{itemize}
    \item $\Omega$ es el espacio muestral (el conjunto de todas las ejecuciones posibles de un algoritmo aleatorizado o juego).
    \item $\mathcal{F}$ es una $\sigma$-álgebra sobre $\Omega$ (conjunto de eventos medibles).
    \item $\mathbb{P}: \mathcal{F} \to [0, 1]$ es la medida de probabilidad.
\end{itemize}
Una \textbf{variable aleatoria} $X$ es una función medible $X: (\Omega, \mathcal{F}) \to (\mathbb{R}, \mathcal{B}(\mathbb{R}))$, donde $\mathcal{B}(\mathbb{R})$ es el álgebra de Borel sobre los reales.

A partir de aquí, estos son los teoremas estructurales y cómo se manifiestan en ICPC.

\section*{1. Funciones Borel-Medibles (LOTUS) y Momentos Superiores}
\textbf{Definición Formal:} Sea $g: \mathbb{R} \to \mathbb{R}$ una función Borel-medible. El teorema conocido como la Ley del Estadístico Inconsciente (LOTUS) establece que el valor esperado de la variable aleatoria $Y = g(X)$ es:
\[
\mathbb{E}[g(X)] = \int_{-\infty}^{\infty} g(x) f_X(x) \, dx
\]
En contextos discretos (competitivas), esto colapsa a $\sum g(x)\mathbb{P}(X=x)$.

\textbf{Insight y Aplicación en ICPC:} \\
La linealidad de la esperanza $\mathbb{E}[X+Y] = \mathbb{E}[X] + \mathbb{E}[Y]$ no requiere independencia, pero $\mathbb{E}[X \cdot Y]$ sí la requiere. El error más común es asumir $\mathbb{E}[X^2] = \mathbb{E}[X]^2$. Cuando un problema te pide el valor esperado del \textbf{cuadrado} o \textbf{cubo} de una cantidad, debes descomponer $X$ en variables indicadoras $I_i \in \{0, 1\}$ y usar LOTUS:
\[
\mathbb{E}[X^2] = \mathbb{E}\left[\left(\sum I_i\right)^2\right] = \sum \mathbb{E}[I_i^2] + 2\sum_{i < j} \mathbb{E}[I_i I_j]
\]
\begin{itemize}
    \item \textbf{Problema Típico:} \textbf{"Inversions after K swaps" (Variante de CF Div. 1)}.
    \item \textbf{Implementación:} Computar $\mathbb{P}(I_i = 1 \cap I_j = 1)$ suele requerir programación dinámica bidimensional u optimizaciones con \textit{Matrix Exponentiation} en un sistema de $O(N^2)$ estados.
\end{itemize}

\section*{2. Filtraciones, Tiempos de Parada e Identidad de Wald}
Respondiendo a tu corrección directa: para la identidad de Wald necesitamos un proceso riguroso.

\textbf{Definición Formal:} Sea $\{X_i\}_{i \geq 1}$ una sucesión de variables aleatorias independientes e idénticamente distribuidas (i.i.d.), donde cada $X_i$ tiene la misma distribución que una variable prototipo $X$. Sea $\{\mathcal{F}_n\}$ su filtración natural $(\mathcal{F}_n = \sigma(X_1, \dots, X_n))$.

Un \textbf{tiempo de parada} (stopping time) es una variable aleatoria discreta $T: \Omega \to \mathbb{N} \cup \{\infty\}$ tal que el evento $\{T = n\} \in \mathcal{F}_n$ (es decir, la decisión de detenerse en el paso $n$ depende solo de la información hasta el paso $n$).

\textbf{Identidad de Wald:} Si $\mathbb{E}[|X|] < \infty$ y $\mathbb{E}[T] < \infty$, entonces la suma aleatoria $S_T = \sum_{i=1}^T X_i$ cumple:
\[
\mathbb{E}[S_T] = \mathbb{E}[T] \cdot \mathbb{E}[X]
\]

\textbf{Insight y Aplicación en ICPC:} \\
Se utiliza para calcular el "costo esperado" o "daño esperado" en juegos donde la duración del juego $T$ es en sí misma aleatoria (por ejemplo, tiras un dado hasta que salga \textbf{6}, y en cada turno ganas una cantidad aleatoria $X_i$ de puntos). En lugar de calcular el DP conjunto de (Puntos, Turnos), la identidad de Wald desacopla el problema: calculas $\mathbb{E}[T]$ (generalmente modelado con Distribución Geométrica) y lo multiplicas por $\mathbb{E}[X]$.
\begin{itemize}
    \item \textbf{Problema Típico:} \textbf{"Collecting Bugs" (POJ 2151)} o problemas de la familia del \textit{Coupon Collector}.
    \item \textbf{Complejidad:} Reduce un problema de estado infinito a dos cálculos independientes de $O(N)$ o incluso $O(1)$ con fórmulas cerradas.
\end{itemize}

\section*{3. Distribuciones Conjugadas y Actualización Bayesiana (Beta-Binomial)}
Este es el núcleo matemático exacto de tu problema \textbf{"Door 1"}.

\textbf{Definición Formal:} La distribución de probabilidad Beta, denotada como $\text{Beta}(\alpha, \beta)$ con $\alpha, \beta > 0$, es una distribución continua en $(0, 1)$ con función de densidad:
\[
f(p; \alpha, \beta) = \frac{p^{\alpha-1}(1-p)^{\beta-1}}{\text{B}(\alpha, \beta)}
\]
donde $\text{B}(\alpha, \beta) = \int_0^1 t^{\alpha-1}(1-t)^{\beta-1} dt$.

\textbf{Teorema (Familias Conjugadas):} La distribución Beta es la \textit{conjugate prior} de la distribución Binomial. Si tenemos una probabilidad $p$ desconocida con distribución previa $p \sim \text{Beta}(\alpha, \beta)$, y realizamos ensayos de Bernoulli observando $s$ éxitos y $f$ fracasos ($n = s+f$), la distribución posterior de $p$ sigue siendo Beta:
\[
p \mid (s, f) \sim \text{Beta}(\alpha + s, \beta + f)
\]

\textbf{Insight y Aplicación en ICPC:} \\
La "Regla de Sucesión de Laplace" es simplemente el caso especial donde el prior es Uniforme en $[0, 1]$. Matemáticamente, la Uniforme $\mathcal{U}(0,1)$ es exactamente $\text{Beta}(1, 1)$.

Tras observar $s$ éxitos y $f$ fracasos, la posterior es $\text{Beta}(1+s, 1+f)$. La esperanza matemática de una distribución $\text{Beta}(\alpha, \beta)$ es $\frac{\alpha}{\alpha+\beta}$.

Por lo tanto, la probabilidad marginal esperada del siguiente éxito es:
\[
\mathbb{E}[p] = \frac{s + 1}{s + f + 2}
\]
\begin{itemize}
    \item \textbf{Implementación en el Problema:} Al saber que el parámetro evoluciona según un proceso markoviano de conteos, evitas integrar en tu DP. Tu estado en "Door 1" solo necesita ser DP(hora, baterias\_actuales, encuentros\_gigante). La probabilidad $p$ del gigante en ese estado está completamente definida algebraicamente por $\frac{\text{encuentros} + 1}{\text{hora} + 2}$.
\end{itemize}

\section*{4. Cadenas de Markov Absorbentes y la Matriz Fundamental}
Cuando abandonamos los DAGs y entramos a grafos con ciclos, el enfoque de DP recursivo puro se rompe por dependencias circulares.

\textbf{Definición Formal:} Sea una Cadena de Markov discreta con estados transitorios y absorbentes. Su matriz de transición en forma canónica por bloques es:
\[
P = \begin{pmatrix} Q & R \\ 0 & I \end{pmatrix}
\]
donde $Q$ representa las transiciones entre estados transitorios, $R$ las transiciones de transitorio a absorbente, e $I$ es la matriz identidad.

\textbf{Teorema:} La matriz fundamental $N$ se define como $N = \sum_{k=0}^{\infty} Q^k = (I - Q)^{-1}$.

Cada entrada $N_{i,j}$ representa el número esperado de visitas al estado transitorio $j$ antes de la absorción, dado que el proceso comenzó en el estado $i$. El tiempo esperado de absorción (número de pasos) partiendo de $i$ es la $i$-ésima entrada del vector $\mathbf{t} = N \mathbf{1}$ (donde $\mathbf{1}$ es un vector columna de unos).

\textbf{Insight y Aplicación en ICPC:} \\
Si un problema te pide el número esperado de pasos para ganar un juego complejo donde puedes "retroceder" o "perder turnos", estás buscando $\mathbf{t}$.
\begin{itemize}
    \item \textbf{Trade-off de Diseño:} Invertir la matriz $(I-Q)$ requiere Eliminación Gaussiana generalizada que corre en $O(V^3)$. Esto restringe el tamaño del grafo de estados a $V \leq 500$. Si la matriz es rala (sparse) o tiene estructura de banda, puedes resolver el sistema lineal iterativamente o usar DP con precisión requerida flotante (ej. iterar hasta que la norma del error sea menor a $10^{-9}$).
    \item \textbf{Problema Típico:} \textbf{"Game of Goose" (Cualquier Regional/WF con tableros circulares y trampas)}.
\end{itemize}

Para alcanzar el nivel de dominio que exige una \textbf{World Finals de la ICPC} o resolver consistentemente los problemas más densos de \textbf{Codeforces Div. 1}, debemos expandir tu arsenal probabilístico con herramientas avanzadas de la teoría de procesos estocásticos y la combinatoria analítica.

A continuación, formalizamos los teoremas avanzados omitidos, detallando su espacio de medida, su comportamiento algebraico y la heurística exacta que te indicará cuándo utilizarlos.

\section*{5. Teorema de Parada Opcional de Doob (Martingalas)}
Cuando un proceso aleatorio posee una propiedad de "juego justo", las técnicas estándar de programación dinámica (DP) o cadenas de Markov colapsan debido a la infinitud del espacio de estados o la complejidad de las transiciones. Aquí es donde las martingalas simplifican el problema radicalmente.

\subsection*{Definición Formal}
Sea $(\Omega, \mathcal{F}, \mathbb{P})$ un espacio de probabilidad equipado con una filtración $\{\mathcal{F}_n\}_{n \geq 0}$. Un proceso estocástico $\{M_n\}_{n \geq 0}$ adaptado a la filtración es una \textbf{martingala} si para todo $n \geq 0$:
\begin{itemize}
    \item $\mathbb{E}[|M_n|] < \infty$
    \item $\mathbb{E}[M_{n+1} \mid \mathcal{F}_n] = M_n$ con probabilidad 1.
\end{itemize}

\subsection*{Teorema de Parada Opcional de Doob (OST)}
Sea $\{M_n\}_{n \geq 0}$ una martingala y $T$ un tiempo de parada respecto a $\{\mathcal{F}_n\}$. Si se cumple al menos una de las siguientes condiciones:
\begin{itemize}
    \item $T$ está acotado casi seguramente ($T \le C$ para alguna constante $C$).
    \item $|M_{n \land T}|$ está acotado uniformemente por una variable integrable.
    \item $\mathbb{E}[T] < \infty$ y existe $K$ tal que $|M_{n+1} - M_n| \le K$ para todo $n < T$.
\end{itemize}
Entonces, el valor esperado de la martingala en el tiempo de parada es igual a su valor esperado inicial:
\[
\mathbb{E}[M_T] = \mathbb{E}[M_0]
\]

\subsection*{Insight y Aplicación en ICPC}
\textbf{¿Cuándo identificarlo?} El problema involucra un juego de apuestas, un \textit{Random Walk} asimétrico en múltiples dimensiones, o la generación de cadenas de caracteres aleatorias (problemas de "tipear un string hasta que aparezca el patrón $S$") y te piden el tiempo esperado de finalización $\mathbb{E}[T]$.
\begin{itemize}
    \item \textbf{El "Truco" de Construcción:} Diseñas una martingala artificial de la forma $M_n = X_n - n \cdot \mathbb{E}[\Delta]$, o basándote en potencias si el proceso es multiplicativo. Al aplicar el OST, resuelves directamente para $\mathbb{E}[T]$ planteando la ecuación $\mathbb{E}[M_T] = \mathbb{E}[M_0]$.
    \item \textbf{Problema de Ejemplo:} \textbf{"K-Consecutive Heads"} o \textbf{"String Matching Games" (Codeforces Div. 1)}. Encontrar el número esperado de lanzamientos de moneda hasta obtener el patrón HTH.
\end{itemize}

\section*{6. Funciones Generatrices de Probabilidad (PGF) y Convolución}
Cuando necesitas calcular no solo la esperanza, sino la distribución de probabilidad exacta de una suma de variables independientes donde los estados tradicionales de DP causarían \textit{Time Limit Exceeded} ($O(N \cdot K)$).

\subsection*{Definición Formal}
Sea $X$ una variable aleatoria discreta que toma valores en los enteros no negativos $\mathbb{N}_0$. La \textbf{Función Generatriz de Probabilidad (PGF)} de $X$ es una función analítica definida como la serie de potencias formal:
\[
G_X(z) = \mathbb{E}[z^X] = \sum_{k=0}^{\infty} \mathbb{P}(X=k) z^k
\]
La serie converge absolutamente para todo $z \in \mathbb{C}$ con $|z| \le 1$.

\subsection*{Propiedades y Teoremas Clave}
\begin{itemize}
    \item \textbf{Convolución:} Si $X$ e $Y$ son variables aleatorias \textbf{independientes}, la PGF de su suma es el producto algebraico de sus PGFs:
    \[
    G_{X+Y}(z) = G_X(z) \cdot G_Y(z)
    \]
    \item \textbf{Recuperación de Momentos:} Los momentos de $X$ se obtienen evaluando las derivadas de $G_X(z)$ en $z=1$:
    \[
    \mathbb{E}[X] = G_X'(1), \quad \mathbb{E}[X(X-1)] = G_X''(1)
    \]
\end{itemize}

\subsection*{Insight y Aplicación en ICPC}
\textbf{¿Cuándo identificarlo?} Te piden calcular la probabilidad de que la suma de $N$ variables independientes (como lanzar $N$ dados deformes con distintas caras) sea exactamente igual a $K$, con $N, K \le 10^5$.
\begin{itemize}
    \item \textbf{El "Path" Mental:} Una DP clásica tardaría $O(N \cdot K)$, lo cual es inviable. Al modelar cada dado como un polinomio $P_i(z) = \sum p_k z^k$, el problema se reduce a calcular el producto de polinomios $\prod_{i=1}^N P_i(z)$.
    \item \textbf{Optimización:} Utilizas un enfoque de \textit{Divide and Conquer} combinado con la \textbf{Transformada Rápida de Fourier (FFT / NTT)} para multiplicar los polinomios, reduciendo la complejidad total a un óptimo:
    \[
    O(K \log K)
    \]
\end{itemize}

\section*{7. Desigualdades de Concentración (Tail Bounds)}
En problemas de algoritmos aleatorizados o estructuras de datos probabilísticas (como \textit{Skip Lists} o \textit{Treaps} modificados), a veces no necesitas el valor exacto de la probabilidad, sino garantizar que la probabilidad de fallo sea infinitesimalmente pequeña.

\subsection*{Teorema (Desigualdad de Chernoff)}
Sean $X_1, X_2, \dots, X_n$ variables aleatorias de Bernoulli independientes (no necesariamente idénticas) tales que $\mathbb{P}(X_i = 1) = p_i$. Sea $X = \sum_{i=1}^n X_i$ y $\mu = \mathbb{E}[X] = \sum p_i$. Para cualquier $\delta > 0$:
\[
\mathbb{P}(X \ge (1+\delta)\mu) \le \left( \frac{e^\delta}{(1+\delta)^{1+\delta}} \right)^\mu
\]

\subsection*{Teorema (Desigualdad de Azuma-Hoeffding)}
Sea $\{M_n\}_{n=0}^m$ una martingala respecto a una filtración, tal que existen constantes $c_n > 0$ que acotan las diferencias casi seguramente: $|M_n - M_{n-1}| \le c_n$. Entonces para cualquier $t > 0$:
\[
\mathbb{P}(M_m - M_0 \ge t) \le \exp\left( -\frac{2t^2}{\sum_{n=1}^m c_n^2} \right)
\]

\subsection*{Insight y Aplicación en ICPC}
\textbf{¿Cuándo identificarlo?} Problemas interactivos o de diseño de algoritmos donde juegas contra un "adversario adaptativo". Te piden diseñar una estrategia que gane con una probabilidad acumulada de $1 - 10^{-9}$.
\begin{itemize}
    \item \textbf{Uso Teórico:} Te permite demostrar por qué un algoritmo codicioso aleatorizado o un muestreo de Monte Carlo con un número fijo de iteraciones (ej. $K = 100$ pruebas) es suficiente para pasar todos los casos de prueba ocultos del juez sin dar \textit{Wrong Answer}.
    \item \textbf{Problema de Ejemplo:} Problemas de geometría computacional donde se realiza un \textit{Randomized Incremental Construction} (como calcular el \textit{Enclosing Circle} mínimo en $O(N)$ esperado).
\end{itemize}

\section*{8. El Principio de Inclusión-Exclusión Probabilístico (Sieve Methods)}
Cuando los eventos cuyas probabilidades deseas calcular están fuertemente interconectados por condiciones de intersección.

\subsection*{Definición Formal}
Sea $(\Omega, \mathcal{F}, \mathbb{P})$ el espacio de medida y $A_1, A_2, \dots, A_n \in \mathcal{F}$ un conjunto de eventos finitos. La medida de la unión de estos eventos está dada por:
\[
\mathbb{P}\left(\bigcup_{i=1}^n A_i\right) = \sum_{i=1}^n \mathbb{P}(A_i) - \sum_{i<j} \mathbb{P}(A_i \cap A_j) + \dots + (-1)^{n-1} \mathbb{P}\left(\bigcap_{i=1}^n A_i\right)
\]

\subsection*{Insight y Aplicación en ICPC}
\textbf{¿Cuándo identificarlo?} Te piden la probabilidad de que una estructura (como un grafo aleatorio o una permutación) cumpla \textbf{al menos una} de varias propiedades malas, o que sea completamente "válida" (ninguna propiedad mala ocurre).
\begin{itemize}
    \item \textbf{El "Path" Mental:} Calcular la probabilidad directa es un problema NP-hard debido a las dependencias. Sin embargo, calcular la probabilidad de la intersección $\mathbb{P}(A_i \cap A_j \cap \dots)$ suele ser trivial si la intersección fija ciertos elementos de la estructura.
    \item \textbf{Integración con Bitmask:} Si $n \le 20$, puedes iterar explícitamente sobre todas las máscaras de bits ($2^n$) para calcular el término del principio de inclusión-exclusión. Si $n$ es grande pero las intersecciones solo dependen del tamaño del subconjunto elegido, se utiliza DP para agrupar los términos por cardinalidad.
    \item \textbf{Problema de Ejemplo:} \textbf{"Derrangements" Probabilísticos} o \textbf{"Counting Graphs" (Latin American Regionals)}.
\end{itemize}

\section*{Hoja de Ruta de Prerrequisitos Avanzados para el Equipo}
Para consolidar esto y evitar errores de precisión o desbordamiento en el código:

\begin{verbatim}
[Teoría de la Medida] ──> [Espacio de Kolmogorov] ──> [Esperanza Condicional Formal]
                                                               │
       ┌────────────────────────┬──────────────────────────────┴────────────────────────┐
       ▼                        ▼                                                       ▼
[Martingalas (OST)]    [Familias Conjugadas (Beta-Binomial)]             [Análisis Polinomial (PGF)]
       │                        │                                                       │
       ▼                        ▼                                                       ▼
Stopping Times discretos     Cálculo analítico de transiciones sin bucles            Multiplicación vía NTT/FFT
\end{verbatim}

\begin{itemize}
    \item \textbf{Aritmética Modular en Probabilidad:} Asegúrense de dominar el cálculo de \textbf{Inversos Modulares} ($\text{mod } 10^9+7$). Nunca operen con \texttt{double} si el problema pide el resultado como $P \cdot Q^{-1} \pmod M$.
    \item \textbf{Propiedad de Markov sin Memoria (Memoryless):} Para variables geométricas discretas, recordar siempre que $\mathbb{P}(X > n+m \mid X > n) = \mathbb{P}(X > m)$. Esto simplifica drásticamente el rediseño de estados en DPs complejas.
\end{itemize}



\end{document}
